\subsection{EXISTENCE THEOREM}

We retain the hypotheses and notation of the preceding number. Recall that if \(\alpha, \beta \in \mathrm{B}\) and if \(\alpha \neq \beta\) , then \(n(\beta, \alpha) \leq 0\) ; moreover, if \(n(\beta, \alpha) = 0\) , then \(n(\alpha, \beta) = 0\) (Chap. VI, §1, no. 1, formula (8)). For any pair \((\alpha, \beta)\) of distinct elements of B, put

\[
x _ {\alpha \beta} = (\mathrm{ad} x _ {\alpha}) ^ {1 - n (\beta , \alpha)} x _ {\beta} y _ {\alpha \beta} = (\mathrm{ad} x _ {- \alpha}) ^ {1 - n (\beta , \alpha)} x _ {- \beta}.
\]

Then \(x_{\alpha \beta} \in \mathfrak{a}_{+}, y_{\alpha \beta} \in \mathfrak{a}_{-}\) . If \(\theta\) denotes the canonical automorphism of \(\mathfrak{a}\) , \(\theta(x_{\alpha \beta}) = y_{\alpha \beta}\) .

Lemma 6. Let \(\alpha, \beta \in \mathrm{B}\) with \(\alpha \neq \beta\) . Then

\[
[ \mathfrak {a} _ {+}, y _ {\alpha \beta} ] = 0 \quad [ \mathfrak {a} _ {-}, x _ {\alpha \beta} ] = 0.
\]

The second formula follows from the first by using the automorphism \(\theta\) . To prove the first, it suffices to show that \([x_{\gamma}, y_{\alpha \beta}] = 0\) for all \(\gamma \in B\) . We distinguish three cases.

Case 1: \(\gamma \neq \alpha\) and \(\gamma \neq \beta\) . In this case, \(x_{\gamma}\) commutes with \(x_{-\alpha}\) and \(x_{-\beta}\) , and hence with \(y_{\alpha \beta}\) .

Case 2: \(\gamma = \beta\) . In this case, \(x_{\gamma}\) commutes with \(x_{-\alpha}\) , so

\[
\begin{array}{r l} & {[ x _ {\gamma}, y _ {\alpha \beta} ] = (\mathrm{ad} x _ {- \alpha}) ^ {1 - n (\beta , \alpha)} [ x _ {\gamma}, x _ {- \beta} ]} \\ & {\qquad = - (\mathrm{ad} x _ {- \alpha}) ^ {1 - n (\beta , \alpha)} h _ {\beta} = - n (\alpha , \beta) (\mathrm{ad} x _ {- \alpha}) ^ {- n (\beta , \alpha)} x _ {- \alpha}.} \end{array}
\]

If \(n(\beta, \alpha) < 0\) , this expression is zero since (ad \(x_{-\alpha}\) ). \(x_{-\alpha} = 0\) . If \(n(\beta, \alpha) = 0\) , then \(n(\alpha, \beta) = 0\) . In both cases, \([x_{\gamma}, y_{\alpha\beta}] = 0\) .

Case 3: \(\gamma = \alpha\) . In the algebra of endomorphisms of \(\mathfrak{a}\) ,

\[
[ - \operatorname{ad} h _ {\alpha}, \operatorname{ad} x _ {- \alpha} ] = 2 \operatorname{ad} x _ {- \alpha}
\]

and {[}ad \(x_{\alpha}\) , ad \(x_{-\alpha}\) {]} = -ad \(h_{\alpha}\) ; thus, by §1, Lemma 1,

\[
[ \operatorname{ad} x _ {\alpha}, (\operatorname{ad} x _ {- \alpha}) ^ {1 - n (\beta , \alpha)} ] = (1 - n (\beta , \alpha)) (\operatorname{ad} x _ {- \alpha}) ^ {- n (\beta , \alpha)} (- \operatorname{ad} h _ {\alpha} - n (\beta , \alpha)).
\]

Consequently,

\[
\begin{array}{c} [ x _ {\gamma}, y _ {\alpha \beta} ] = [ \text {ad} x _ {\alpha}, (\text {ad} x _ {- \alpha}) ^ {1 - n (\beta , \alpha)} ] x _ {- \beta} + (\text {ad} x _ {- \alpha}) ^ {1 - n (\beta , \alpha)} (\text {ad} x _ {\alpha}) x _ {- \beta} \\ = - (1 - n (\beta , \alpha)) (\text {ad} x _ {- \alpha}) ^ {- n (\beta , \alpha)} (\text {ad} h _ {\alpha} + n (\beta , \alpha)) x _ {- \beta} \\ + (\text {ad} x _ {- \alpha}) ^ {1 - n (\beta , \alpha)} (\text {ad} x _ {\alpha}) x _ {- \beta}. \end{array}
\]

Now \([h_{\alpha},x_{-\beta}] + n(\beta ,\alpha)x_{-\beta} = 0\) and \([x_{\alpha},x_{-\beta}] = 0\) , so \([x_{\gamma},y_{\alpha \beta}] = 0\) .

Lemma 7. The ideal \(\mathfrak{n}\) of \(\mathfrak{a}_{+}\) generated by the \(x_{\alpha \beta} (\alpha, \beta \in \mathrm{B}, \alpha \neq \beta)\) is an ideal of \(\mathfrak{a}\) . The ideal of \(\mathfrak{a}_{-}\) generated by the \(y_{\alpha \beta} (\alpha, \beta \in \mathrm{B}, \alpha \neq \beta)\) is an ideal of \(\mathfrak{a}\) and is equal to \(\theta(\mathfrak{n})\) .

Let \(\mathfrak{n}' = \sum_{\alpha, \beta \in \mathrm{B}, \alpha \neq \beta} kx_{\alpha \beta}\) . Since each \(x_{\alpha \beta}\) is homogeneous in \(\mathfrak{a}\) , \([\mathfrak{a}^0, \mathfrak{n}'] \subset \mathfrak{n}'\) (Lemma 5 and Prop. 2). Let U (resp. V) be the enveloping algebra of \(\mathfrak{a}\) (resp. \(\mathfrak{a}_+\) ), and \(\sigma\) the representation of U on \(\mathfrak{a}\) defined by the adjoint representation of \(\mathfrak{a}\) . The ideal of \(\mathfrak{a}\) generated by \(\mathfrak{n}'\) is \(\sigma(\mathrm{U})\mathfrak{n}'\) . Now \(\mathfrak{a} = \mathfrak{a}_+ + \mathfrak{a}^0 + \mathfrak{a}_-\) (Prop. 2), \(\sigma(\mathfrak{a}_-)\mathfrak{n}' = 0\) (Lemma 6), and \(\sigma(\mathfrak{a}^0)\mathfrak{n}' \subset \mathfrak{n}'\) by the preceding. By the Poincaré-Birkhoff-Witt theorem, \(\sigma(\mathrm{U})\mathfrak{n}' = \sigma(\mathrm{V})\mathfrak{n}'\) , which proves the first assertion of the lemma. It follows that the ideal of \(\theta(\mathfrak{a}_+) = \mathfrak{a}_-\) generated by the \(\theta(x_{\alpha \beta}) = y_{\alpha \beta}\) ( \(\alpha, \beta \in \mathrm{B}, \alpha \neq \beta\) ) is the ideal \(\theta(\mathfrak{n})\) of \(\mathfrak{a}\) . Q.E.D.

The ideal \(\mathfrak{n} + \theta(\mathfrak{n})\) of a is graded since it is generated by homogeneous elements. Consequently, the Lie algebra \(\mathfrak{a}/(\mathfrak{n}+\theta(\mathfrak{n}))\) is a Q-graded Lie algebra; in the remainder of this paragraph, it is denoted by \(g_{B}\) , or simply by g. By Prop. 2, if \(g^{\mu} \neq 0\) then \(\mu \in Q_{+}\) , or \(\mu \in Q_{-}\) , or \(\mu = 0\) . Denote by \(X_{\alpha}\) (resp. \(H_{\alpha}, X_{-\alpha}\) ) the canonical image of \(x_{\alpha}\) (resp. \(h_{\alpha}, x_{-\alpha}\) ) in g. In view of the definition of a, n and \(\theta(\mathfrak{n})\) , it follows that g is the Lie algebra defined by the generating family \(((X_{\alpha}, H_{\alpha}, X_{-\alpha}))_{\alpha \in B}\) and the relations

\[
[ H _ {\alpha}, H _ {\beta} ] = 0\tag{16}
\]

\[
[ H _ {\alpha}, X _ {\beta} ] - n (\beta , \alpha) X _ {\beta} = 0\tag{17}
\]

\[
[ H _ {\alpha}, X _ {- \beta} ] + n (\beta , \alpha) X _ {- \beta} = 0\tag{18}
\]

\[
[ X _ {\alpha}, X _ {- \alpha} ] + H _ {\alpha} = 0\tag{19}
\]

\[
[ X _ {\alpha}, X _ {- \beta} ] = 0 (\alpha \neq \beta)\tag{20}
\]

\[
(\operatorname{ad} X _ {\alpha}) ^ {1 - n (\beta , \alpha)} X _ {\beta} = 0 \quad (\alpha \neq \beta)\tag{21}
\]

\[
(\operatorname{ad} X _ {- \alpha}) ^ {1 - n (\beta , \alpha)} X _ {- \beta} = 0 \quad (\alpha \neq \beta).\tag{22}
\]

Let \(z \in \mathfrak{g}\) and \(\mu \in \mathrm{Q}\) . Then \(z \in \mathfrak{g}^{\mu}\) if and only if \([H_{\alpha}, z] = \langle \mu, \alpha^{\vee} \rangle z\) for all \(\alpha \in \mathrm{B}\) . This follows from Lemma 5.

Since \(\mathfrak{a}^0\cap (\mathfrak{n} + \theta (\mathfrak{n})) = 0\) , the canonical map from \(\mathfrak{a}^0\) to \(\mathfrak{g}^0\) is an isomorphism. Consequently, \((H_{\alpha})_{\alpha \in \mathrm{B}}\) is a basis of the vector space \(\mathfrak{g}^0\) . The commutative subalgebra \(\mathfrak{g}^0\) of \(\mathfrak{g}_{\mathrm{B}}\) will be denoted by \(\mathfrak{h}_{\mathrm{B}}\) or simply by \(\mathfrak{h}\) . There exists a unique isomorphism \(\mu \mapsto \mu_{\mathrm{B}}\) from V to \(\mathfrak{h}^*\) such that \(\langle \mu_{\mathrm{B}},H_{\alpha}\rangle = \langle \mu ,\alpha^{\vee}\rangle\) for all \(\mu \in \mathrm{V}\) and all \(\alpha \in \mathrm{B}\) .

The involutive automorphism \(\theta\) of a defines by passage to the quotient an involutive automorphism of g that will also be denoted by \(\theta\) . We have \(\theta(X_{\alpha}) = X_{-\alpha}\) for \(\alpha \in B \cup (-B)\) , and \(\theta(H_{\alpha}) = -H_{\alpha}\) .

THEOREM 1. Let R be a reduced root system, B a basis of R. Let g be the Lie algebra defined by the generating family \(((X_{\alpha}, H_{\alpha}, X_{-\alpha}))_{\alpha \in \mathrm{B}}\) and the relations (16) to (22). Let \(h = \sum_{\alpha \in B} kH_{\alpha}\) . Then \((\mathfrak{g}, \mathfrak{h})\) is a split semi-simple Lie algebra. The isomorphism \(\mu \mapsto \mu_{B}\) from V to \(h^{*}\) maps R to the root system of \((\mathfrak{g}, \mathfrak{h})\) . For all \(\mu \in R\) , \(g^{\mu}\) is the eigenspace relative to the root \(\mu\) .

The proof follows that of Lemmas 8, 9, 10, 11.

Lemma 8. Let \(\alpha \in \mathrm{B} \cup (-\mathrm{B})\) . Then \(\operatorname{ad} X_{\alpha}\) is locally nilpotent. \footnote{An endomorphism \(u\) of a vector space \(V\) is called locally nilpotent (or almost nilpotent) if, for every \(v \in V\), there exists a positive integer \(n\) such that \(u^n(v) = 0\) (cf. Commutative Algebra, Chap. IV, §1, no. 4, Def. 2). Then \(\exp(u)\), or \(e^u\), is defined by the formula \(e^u(v) = \sum_{n \geq 0} (1/n!)u^n(v)\) for all \(v \in V\).}

Assume that \(\alpha \in \mathrm{B}\) . Let \(\mathfrak{g}'\) be the set of \(z \in \mathfrak{g}\) such that \((\operatorname{ad} X_{\alpha})^p z = 0\) for sufficiently large \(p\) . Since \(\operatorname{ad} X_{\alpha}\) is a derivation of \(\mathfrak{g}\) , \(\mathfrak{g}'\) is a subalgebra of \(\mathfrak{g}\) . By (21), \(X_{\beta} \in \mathfrak{g}'\) for all \(\beta \in \mathrm{B}\) . By (17), (19), (20), \(H_{\beta} \in \mathfrak{g}'\) and \(X_{-\beta} \in \mathfrak{g}'\) for all \(\beta \in \mathrm{B}\) . Hence \(\mathfrak{g}' = \mathfrak{g}\) and \(\operatorname{ad} X_{\alpha}\) is locally nilpotent. Since \(\operatorname{ad} X_{-\alpha} = \theta (\operatorname{ad} X_{\alpha})\theta^{-1}\) , we see that \(\operatorname{ad} X_{-\alpha}\) is locally nilpotent.

We shall see that g is finite dimensional, so that \(ad X_{\alpha}\) is actually nilpotent.

Lemma 9. Let \(\mu, \nu \in \mathbf{Q}\) and \(w \in \mathrm{W}(\mathrm{R})\) be such that \(w\mu = \nu\) . There exists an automorphism of \(\mathfrak{g}\) that takes \(\mathfrak{g}^{\mu}\) to \(\mathfrak{g}^{\nu}\) .

For all \(\alpha \in \mathrm{B}\) , let \(s_{\alpha}\) be the reflection in V defined by \(\alpha\) . Since \(\mathrm{W}(\mathrm{R})\) is generated by the \(s_{\alpha}\) (Chap. VI, §1, no. 5, Remark 1), it suffices to prove the lemma when \(w = s_{\alpha}\) . In view of Lemma 8, we can define

\[
\theta_ {\alpha} = e ^ {\mathrm{ad} X _ {\alpha}} e ^ {\mathrm{ad} X _ {- \alpha}} e ^ {\mathrm{ad} X _ {\alpha}}.
\]

It is verified as in Chap. I, §6, no. 8, that \(\theta_{\alpha}\) is an automorphism of \(\mathfrak{g}\) . We have

\[
\begin{array}{l} \theta_ {\alpha} (H _ {\alpha}) = e ^ {\mathrm{ad} X _ {\alpha}} e ^ {\mathrm{ad} X _ {- \alpha}} (H _ {\beta} - n (\alpha , \beta) X _ {\alpha}) \\ \qquad = e ^ {\mathrm{ad} X _ {\alpha}} \bigg (H _ {\beta} - n (\alpha , \beta) X _ {\alpha} + n (\alpha , \beta) X _ {- \alpha} - n (\alpha , \beta) H _ {\alpha} - \frac {n (\alpha , \beta)}{2} 2 X _ {- \alpha} \bigg) \\ \qquad = e ^ {\mathrm{ad} X _ {\alpha}} \left(H _ {\beta} - n (\alpha , \beta) H _ {\alpha} - n (\alpha , \beta) X _ {\alpha}\right) \\ \qquad = H _ {\beta} - n (\alpha , \beta) H _ {\alpha} - n (\alpha , \beta) X _ {\alpha} - n (\alpha , \beta) X _ {\alpha} - n (\alpha , \beta) (- 2 X _ {\alpha}) \\ \qquad = H _ {\beta} - n (\alpha , \beta) H _ {\alpha}. \end{array}
\]

If \(z\in \mathfrak{g}^{\mu}\)

\[
\begin{array}{r l} & {[ H _ {\beta}, \theta_ {\alpha} ^ {- 1} z ] = \theta_ {\alpha} ^ {- 1} [ H _ {\beta} - n (\alpha , \beta) H _ {\alpha}, z ]} \\ & {\qquad = \theta_ {\alpha} ^ {- 1} (\langle \mu , \beta^ {\vee} \rangle z - n (\alpha , \beta) \langle \mu , \alpha^ {\vee} \rangle z)} \\ & {\qquad = \langle \mu - \langle \alpha^ {\vee}, \mu \rangle \alpha , \beta^ {\vee} \rangle \theta_ {\alpha} ^ {- 1} z = \langle s _ {\alpha} \mu , \beta^ {\vee} \rangle \theta_ {\alpha} ^ {- 1} z,} \end{array}
\]

so \(\theta_{\alpha}^{-1}z\in\mathfrak{g}^{s_{\alpha}\mu}\) . This shows that \(\theta_{\alpha}^{-1}\mathfrak{g}^{\mu}\subset\mathfrak{g}^{s_{\alpha}\mu}\) . Since \(\theta_{\alpha}\) is an automorphism and since this inclusion holds for all \(\mu\in Q\) , we see that \(\theta_{\alpha}^{-1}\mathfrak{g}^{\mu}=\mathfrak{g}^{s_{\alpha}\mu}\) , which proves the lemma.

Lemma 10. Let \(\mu\in Q\) , and assume that \(\mu\) is not a multiple of a root. There exists \(w\in\mathrm{W}(\mathrm{R})\) such that certain of the coordinates of \(w\mu\) with respect to the basis B are \(> 0\) and certain of them are \(< 0\).

Let \(V_{R}\) be the vector space \(Q \otimes_{Z} R\) , in which R is a root system. By the assumption, there exists \(f \in V_{R}^{*}\) such that \(\langle f, \alpha \rangle \neq 0\) for all \(\alpha \in R\) , and \(\langle f, \mu \rangle = 0\) . There exists a chamber C of \(R^{\vee}\) such that \(f \in C\) . By Chap. VI, §1, no. 5, Th. 2 (i), there exists \(w \in W(R)\) such that wf belongs to the chamber associated to B, in other words such that \(\langle wf, \alpha \rangle > 0\) for all \(\alpha \in B\) . Write \(w\mu = \sum_{\alpha \in B} t_{\alpha}\alpha\) . Then

\[
0 = \langle f, \mu \rangle = \langle w f, w \mu \rangle = \sum_ {\alpha \in B} t _ {\alpha} \langle w f, \alpha \rangle ,
\]

which proves that certain \(t_{\alpha}\) are \(> 0\) and others are \(< 0\) .

Lemma 11. Let \(\mu \in \mathrm{Q}\) . If \(\mu \notin \mathrm{R} \cup \{0\}\) , then \(\mathfrak{g}^{\mu} = 0\) . If \(\mu \in \mathrm{R}\) , then \(\dim \mathfrak{g}^{\mu} = 1\) . 1) If \(\mu\) is not a multiple of an element of \(\mathrm{R}\) , there exists \(w \in \mathrm{W}\) such that \(w\mu \notin \mathrm{Q}_{+} \cup \mathrm{Q}_{-}\) (Lemma 10), so \(\mathfrak{a}^{w\mu} = 0\) , \(\mathfrak{g}^{w\mu} = 0\) , and hence \(\mathfrak{g}^{\mu} = 0\) (Lemma 9).

\begin{enumerate}
\def\labelenumi{\arabic{enumi})}
\setcounter{enumi}{1}
\item
  Let \(\alpha \in \mathrm{B}\) and let \(m\) be an integer. Since \(\mathfrak{a}_{+}\) is a free Lie algebra with basic family \((x_{\alpha})_{\alpha \in \mathrm{B}}\) , we have \(\dim \mathfrak{a}^{\alpha} = 1\) and \(\mathfrak{a}^{m\alpha} = 0\) for \(m > 1\) (Chap. II, §2, no. 6, Prop. 4). Hence \(\dim \mathfrak{g}^{\alpha} \leq 1\) and \(\mathfrak{g}^{m\alpha} = 0\) for \(m > 1\) . We cannot have \(\mathfrak{g}^{\alpha} = 0\) , as this would imply that \(x_{\alpha} \in \mathfrak{n} + \theta \mathfrak{n}\) , and hence that \(\mathfrak{n} + \theta \mathfrak{n}\) contains \(h_{\alpha} = -[x_{\alpha}, x_{-\alpha}]\) , whereas \(\mathfrak{a}^0 \cap (\mathfrak{n} + \theta \mathfrak{n}) = 0\) . Consequently, \(\dim \mathfrak{g}^{\alpha} = 1\) .
\item
  If \(\mu \in \mathrm{R}\) , there exists \(w \in \mathrm{W}(\mathrm{R})\) such that \(w(\mu) \in \mathrm{B}\) (Chap. VI, §1, no. 5, Prop. 15), so \(\dim \mathfrak{g}^{\mu} = \dim \mathfrak{g}^{w\mu} = 1\) . Moreover, if \(n\) is an integer \(>1\) then \(\mathfrak{g}^{nw(\mu)} = 0\) and so \(\mathfrak{g}^{n\mu} = 0\) .
\end{enumerate}

Proof of Theorem 1.

Since \(\dim \mathfrak{g}^0 = \operatorname{CardB}\) , it follows from Lemma 11 that \(\mathfrak{g}\) is of finite dimension equal to \(\operatorname{CardB} + \operatorname{CardR}\) . We show that \(\mathfrak{g}\) is semi-simple. Let \(\mathfrak{k}\) be a commutative ideal of \(\mathfrak{g}\) . Since \(\mathfrak{k}\) is stable under \(\operatorname{ad}(\mathfrak{h})\) , \(\mathfrak{k} = (\mathfrak{k} \cap \mathfrak{h}) + \sum_{\mu \in \mathrm{R}} (\mathfrak{k} \cap \mathfrak{g}^\mu)\) .

It is clear that, for all \(\alpha \in \mathrm{B}\) , \(\mathfrak{g}^{\alpha} + \mathfrak{g}^{-\alpha} + kH_{\alpha}\) is isomorphic to \(\mathfrak{sl}(2,k)\) . In view of Lemma 9, for all \(\mu \in \mathrm{R}\) , \(\mathfrak{g}^{\mu}\) is contained in a subalgebra of \(\mathfrak{g}\) isomorphic to \(\mathfrak{sl}(2,k)\) ; consequently, \(\mathfrak{k} \cap \mathfrak{g}^{\mu} = 0\) , so \(\mathfrak{k} \subset \mathfrak{h}\) ; hence

\[
[ \mathfrak {k}, \mathfrak {g} ^ {\mu} ] \subset \mathfrak {k} \cap \mathfrak {g} ^ {\mu} = 0,
\]

so \(\mu_{\mathrm{B}}(\mathfrak{k}) = 0\) for all \(\mu \in \mathrm{R}\) . It follows that \(\mathfrak{k} = 0\) , which proves that \(\mathfrak{g}\) is semi-simple.

Let \(\mu \in \mathrm{R}\) . There exists \(\alpha \in \mathrm{B}\) such that \(\langle \mu, \alpha^{\vee} \rangle \neq 0\) , and \((\operatorname{ad} H_{\alpha})|_{\mathfrak{g}^{\mu}}\) is then a non-zero homothety. Consequently, \(\mathfrak{h}\) is equal to its own normalizer in \(\mathfrak{g}\) , and hence is a Cartan subalgebra of \(\mathfrak{g}\) . For all \(u \in \mathfrak{h}\) , \(\operatorname{ad} u\) is diagonalizable, so \((\mathfrak{g}, \mathfrak{h})\) is a split semi-simple Lie algebra.

For all \(\mu \in \mathrm{R}\) , it is clear that \(\mu_{\mathrm{B}}\) is a root of \((\mathfrak{g},\mathfrak{h})\) and that \(\mathfrak{g}^{\mu}\) is the corresponding eigenspace. The number of roots of \((\mathfrak{g},\mathfrak{h})\) is \(\dim \mathfrak{g} - \dim \mathfrak{h} = \operatorname{Card}\mathrm{R}\) . Hence, the map \(\mu \mapsto \mu_{\mathrm{B}}\) from V to \(\mathfrak{h}^*\) maps R to the root system of \((\mathfrak{g},\mathfrak{h})\) .

\subsection{UNIQUENESS THEOREM}

PROPOSITION 4. Let \((\mathfrak{g},\mathfrak{h},\mathrm{B},(X_{\alpha})_{\alpha\in\mathrm{B}})\) be a framed semi-simple Lie algebra. Let \((n(\alpha,\beta))_{\alpha,\beta\in\mathrm{B}}\) and \((X_{\alpha})_{\alpha\in\mathrm{B}\cup(-\mathrm{B})}\) be the corresponding Cartan matrix and generating family.

\begin{enumerate}
\def\labelenumi{(\roman{enumi})}
\item
  The family \(((X_{\alpha}, H_{\alpha}, X_{-\alpha}))_{\alpha \in \mathrm{B}}\) and the relations (16) to (22) of no. 3 constitute a presentation of \(\mathfrak{g}\) .
\item
  The family \((X_{\alpha})_{\alpha \in \mathrm{B}}\) and the relations (21) of no. 3 constitute a presentation of the subalgebra of \(\mathfrak{g}\) generated by \((X_{\alpha})_{\alpha \in \mathrm{B}}\) .
\end{enumerate}

Let R be the root system of \((\mathfrak{g},\mathfrak{h})\) . Applying to R and B the constructions of nos. 2 and 3, we obtain objects that we shall denote by \(a', g', X'_{\alpha}, H'_{\alpha}, \ldots\) instead of \(a, g, X_{\alpha}, H_{\alpha}, \ldots\)

There exists a homomorphism \(\varphi\) from the Lie algebra \(\mathfrak{g}'\) to the Lie algebra \(\mathfrak{g}\) such that \(\varphi(X_{\alpha}') = X_{\alpha}\) , \(\varphi(H_{\alpha}') = H_{\alpha}\) , \(\varphi(X_{-\alpha}') = X_{-\alpha}\) for all \(\alpha \in B\) (Prop. 1). Since \(\dim \mathfrak{g}' = \operatorname{Card} R + \operatorname{Card} B = \dim \mathfrak{g}\) , \(\varphi\) is bijective. This proves (i).

The subalgebra of \(\mathfrak{g}' = \mathfrak{a}' / (\mathfrak{n}'\oplus \theta'\mathfrak{n}') = (\mathfrak{a}_+^\prime \oplus \mathfrak{a}^{\prime 0}\oplus \mathfrak{a}_-^\prime) / (\mathfrak{n}'\oplus \theta'\mathfrak{n}')\) generated by \((X_{\alpha}^{\prime})_{\alpha \in \mathrm{B}}\) can be identified with \(\mathfrak{a}_+^\prime /\mathfrak{n}'\) . In view of Prop. 3 and the definition of \(\mathfrak{n}'\) , this proves (ii).

COROLLARY. Every framed semi-simple Lie algebra is obtained from a framed semi-simple \(\mathbf{Q}\) -Lie algebra by extension of scalars from \(\mathbf{Q}\) to \(k\) .

This follows immediately from the proposition.

THEOREM 2. Let \((\mathfrak{g},\mathfrak{h},\mathrm{B},(X_{\alpha})_{\alpha \in \mathrm{B}})\) and \((\mathfrak{g}',\mathfrak{h}',\mathrm{B}',(X'_{\alpha})_{\alpha \in \mathrm{B}'})\) be framed semi-simple Lie algebras, let R and R' be the root systems of \((\mathfrak{g},\mathfrak{h})\) and \((\mathfrak{g}',\mathfrak{h}')\) , let \((n(\alpha ,\beta))_{\alpha ,\beta \in \mathrm{B}}\) (resp. \((n^{\prime}(\alpha ,\beta))_{\alpha ,\beta \in \mathrm{B}'}\) ) be the Cartan matrix of R (resp. R') relative to B (resp. B'), and let \(\Delta\) (resp. \(\Delta'\) ) be the Dynkin graph of R (resp. R') relative to B (resp. B').

\begin{enumerate}
\def\labelenumi{(\roman{enumi})}
\item
  If \(\varphi\) is an isomorphism from \(\mathfrak{h}^*\) to \(\mathfrak{h}'^*\) such that \(\varphi(\mathrm{R}) = \mathrm{R}'\) and \(\varphi(\mathrm{B}) = \mathrm{B}'\) , there exists a unique isomorphism \(\psi\) from \((\mathfrak{g}, \mathfrak{h}, \mathrm{B}, (X_{\alpha})_{\alpha \in \mathrm{B}})\) to \((\mathfrak{g}', \mathfrak{h}', \mathrm{B}', (X_{\alpha}')_{\alpha \in \mathrm{B}'})\) such that \(\psi|_{\mathfrak{h}} = {}^t\varphi^{-1}\) .
\item
  If \(f\) is a bijection from \(\mathrm{B}\) to \(\mathrm{B}'\) such that \(n'(f(\alpha), f(\beta)) = n(\alpha, \beta)\) for all \(\alpha, \beta \in \mathrm{B}\) , there exists an isomorphism from \((\mathfrak{g}, \mathfrak{h}, \mathrm{B}, (X_{\alpha})_{\alpha \in \mathrm{B}})\) to \((\mathfrak{g}', \mathfrak{h}', \mathrm{B}', (X_{\alpha}')_{\alpha \in \mathrm{B}'}).\)
\item
  If there exists an isomorphism from \(\Delta\) to \(\Delta'\) , there exists an isomorphism from \((\mathfrak{g},\mathfrak{h},\mathrm{B},(X_{\alpha})_{\alpha \in \mathrm{B}})\) to \((\mathfrak{g}',\mathfrak{h}',\mathrm{B}',(X_{\alpha}')_{\alpha \in \mathrm{B}'})\) .
\end{enumerate}

This follows immediately from Prop. 4 (i) (making use of Chap. VI, §4, no. 2, Prop. 1 for part (iii)).

Scholium. To any splittable semi-simple Lie algebra g is associated a Dynkin graph, which determines g up to isomorphism (Th. 2 (iii)). This graph is non-empty and connected if and only if g is simple (§3, no. 2, Cor. 1 of Prop. 6). By Th. 1 of no. 3, and Chap. VI, §4, no. 2, Th. 3, the splittable simple Lie algebras are the algebras of type \(A_{l}\) ( \(l \geq 1\) ), \(B_{l}\) ( \(l \geq 2\) ), \(C_{l}\) ( \(l \geq 3\) ), \(D_{l}\) ( \(l \geq 4\) ), \(E_{6}\) , \(E_{7}\) , \(E_{8}\) , \(F_{4}\) , \(G_{2}\) . No two algebras in this list are isomorphic.

PROPOSITION 5. Let \((\mathfrak{g},\mathfrak{h},\mathrm{B},(X_{\alpha})_{\alpha \in \mathrm{B}})\) be a framed semi-simple Lie algebra, and \((X_{\alpha})_{\alpha \in \mathrm{B}\cup (-\mathrm{B})}\) the corresponding generating family. There exists a unique automorphism \(\theta\) of \(\mathfrak{g}\) such that \(\theta (X_{\alpha}) = X_{-\alpha}\) for all \(\alpha \in \mathrm{B}\cup (-\mathrm{B})\) . We have \(\theta^2 = \mathrm{Id}_{\mathfrak{g}}\) , and \(\theta (h) = -h\) for all \(h\in \mathfrak{h}\) .

The uniqueness is clear since \((X_{\alpha})_{\alpha\in\mathrm{B}\cup(-\mathrm{B})}\) generates the Lie algebra g. In view of Prop. 4, the existence of \(\theta\) follows from what we said in no. 3 before Th. 1.

COROLLARY. Let \((\mathfrak{g},\mathfrak{h})\) be a split semi-simple Lie algebra. Then \((\mathfrak{g},\mathfrak{h})\) possesses a Chevalley system (§2, no. 4, Def. 3).

Let R be the root system of \((\mathfrak{g},\mathfrak{h})\) . For all \(\alpha\in R\) , let \(X_{\alpha}\) be a non-zero element of \(g^{\alpha}\) . Assume that the \(X_{\alpha}\) are chosen so that \([X_{\alpha},X_{-\alpha}]=-H_{\alpha}\) for all \(\alpha\in R\) (§2, no. 4, Lemma 2). Let B be a basis of R and \(\theta\) the automorphism of g such that \(\theta(X_{\alpha})=X_{-\alpha}\) for all \(\alpha\in\mathrm{B}\cup(-\mathrm{B})\) . We have \(\theta|_{\mathfrak{h}}=-\mathrm{Id}_{\mathfrak{h}}\) . Hence, for all \(\alpha\in R\) there exists \(t_{\alpha}\in k^{*}\) such that \(\theta X_{\alpha}=t_{\alpha}X_{-\alpha}\) . We have

\[
\begin{array}{c} t _ {\alpha} t _ {- \alpha} H _ {\alpha} = [ t _ {\alpha} X _ {- \alpha}, t _ {- \alpha} X _ {\alpha} ] = [ \theta X _ {\alpha}, \theta X _ {- \alpha} ] = \theta ([ X _ {\alpha}, X _ {- \alpha} ]) \\ = \theta (- H _ {\alpha}) = H _ {\alpha} \end{array}
\]

so \(t_{\alpha}t_{-\alpha} = 1\) for all \(\alpha \in \mathrm{R}\) . Introduce the \(\mathrm{N}_{\alpha \beta}\) as in §2, no. 4. If \(\alpha, \beta, \alpha + \beta \in \mathrm{R}\) ,

\[
\begin{array}{r} \mathrm{N} _ {- \alpha , - \beta} t _ {\alpha} t _ {\beta} X _ {- \alpha - \beta} = t _ {\alpha} t _ {\beta} [ X _ {- \alpha}, X _ {- \beta} ] = [ \theta X _ {\alpha}, \theta X _ {\beta} ] = \theta ([ X _ {\alpha}, X _ {\beta} ]) \\ = \mathrm{N} _ {\alpha \beta} \theta X _ {\alpha + \beta} = \mathrm{N} _ {\alpha \beta} t _ {\alpha + \beta} X _ {- \alpha - \beta} \end{array}
\]

so, in view of §2, no. 4, Lemma 4,

\[
(q + 1) ^ {2} t _ {\alpha} t _ {\beta} = \mathrm{N} _ {\alpha \beta} ^ {2} t _ {\alpha + \beta}
\]

where \(q\) is an integer. It follows that if \(t_{\alpha}\) and \(t_{\beta}\) are squares in \(k^*\) , so is \(t_{\alpha + \beta}\) . Since \(t_{\alpha} = 1\) for all \(\alpha \in \mathrm{B}\) , Prop. 19 of Chap. VI, §1, no. 6, proves that \(t_{\alpha}\) is a square for all \(\alpha \in \mathrm{R}\) . Choose, for all \(\alpha \in \mathrm{R}\) , a \(u_{\alpha} \in k\) such that \(u_{\alpha}^{2} = t_{\alpha}\) . This choice can be made so that \(u_{\alpha}u_{-\alpha} = 1\) for all \(\alpha \in \mathrm{R}\) . Put \(X_{\alpha}' = u_{\alpha}^{-1}X_{\alpha}\) . Then, for all \(\alpha \in \mathrm{R}\) ,

\[
X _ {\alpha} ^ {\prime} \in \mathfrak {g} ^ {\alpha}, [ X _ {\alpha} ^ {\prime}, X _ {- \alpha} ^ {\prime} ] = [ X _ {\alpha}, X _ {- \alpha} ] = - H _ {\alpha},
\]

and \(\theta(X_{\alpha}^{\prime}) = \theta(u_{\alpha}^{-1}X_{\alpha}) = u_{\alpha}^{-1}t_{\alpha}X_{-\alpha} = u_{\alpha}X_{-\alpha} = u_{\alpha}u_{-\alpha}X_{-\alpha}^{\prime} = X_{-\alpha}^{\prime}\) , so that \((X_{\alpha}^{\prime})_{\alpha \in \mathrm{R}}\) is a Chevalley system of \((\mathfrak{g},\mathfrak{h})\) .

\section{AUTOMORPHISMS OF A SEMI-SIMPLE LIE ALGEBRA}

In this paragraph, g denotes a semi-simple Lie algebra.

\subsection{AUTOMORPHISMS OF A FRAMED SEMI-SIMPLE LIE ALGEBRA}

Recall (Chap. VII, §3, no. 1) that \(\operatorname{Aut}(\mathfrak{g})\) denotes the group of automorphisms of \(\mathfrak{g}\) . If \(\mathfrak{h}\) is a Cartan subalgebra of \(\mathfrak{g}\) , we denote by \(\operatorname{Aut}(\mathfrak{g},\mathfrak{h})\) the group of automorphisms of \(\mathfrak{g}\) that leave \(\mathfrak{h}\) stable. Assume that \(\mathfrak{h}\) is splitting, and let R be the root system of \((\mathfrak{g},\mathfrak{h})\) . If \(s\in\operatorname{Aut}(\mathfrak{g},\mathfrak{h})\) , the contragredient map of \(s|_{\mathfrak{h}}\) is an element of A(R) (the group of automorphisms of R) which we shall denote by \(\varepsilon(s)\) in this paragraph. Thus

\[
\varepsilon : \operatorname{Aut} (\mathfrak {g}, \mathfrak {h}) \to \mathrm{A} (\mathrm{R})
\]

is a homomorphism of groups.

For any root system R and any basis B of R, we denote by \(\operatorname{Aut}(R,B)\) the group of automorphisms of R that leave B stable. Recall (Chap. VI, §1, no. 5, Prop. 16 and §4, no. 2, Cor. of Prop. 1) that A(R) is the semi-direct product of \(\operatorname{Aut}(R,B)\) and W(R), and that A(R)/W(R) is canonically isomorphic to the group of automorphisms of the Dynkin graph of R.

PROPOSITION 1. Let \((\mathfrak{g},\mathfrak{h},\mathrm{B},(X_{\alpha})_{\alpha \in \mathrm{B}})\) be a framed semi-simple Lie algebra, and R the root system of \((\mathfrak{g},\mathfrak{h})\) . Let G be the set of \(s\in \operatorname {Aut}(\mathfrak{g},\mathfrak{h})\) that leave B stable, and such that \(s(X_{\alpha}) = X_{\varepsilon (s)\alpha}\) for all \(\alpha \in \mathrm{B}\) (in other words the set of automorphisms of \((\mathfrak{g},\mathfrak{h},\mathrm{B},(X_{\alpha})_{\alpha \in \mathrm{B}})\) . Then the restriction of \(\varepsilon\) to \(\mathrm{G}\) is an isomorphism from \(\mathrm{G}\) to \(\operatorname {Aut}(\mathrm{R},\mathrm{B})\) .

If \(s \in G\) , it is clear that \(\varepsilon(s) \in \operatorname{Aut}(\mathrm{R}, \mathrm{B})\) . On the other hand, the map

\[
\varepsilon|_{\mathrm{G}}: \mathrm{G} \rightarrow \operatorname{Aut} (\mathrm{R}, \mathrm{B})
\]

is bijective by Th. 2 of §4, no. 4.

\subsection{AUTOMORPHISMS OF A SPLIT SEMI-SIMPLE LIE ALGEBRA}

Let E be a commutative group, and \(A = \bigoplus_{\gamma \in E} A^{\gamma}\) an E-graded algebra. For any homomorphism \(\varphi\) from the group E to the multiplicative group \(k^{*}\) , let \(f(\varphi)\) be the k-linear map from A to A whose restriction to each \(A^{\gamma}\) is the homothety with ratio \(\varphi(\gamma)\) ; it is clear that \(f(\varphi)\) is an automorphism of the graded algebra A, and that f is a homomorphism from the group \(\operatorname{Hom}(E, k^{*})\) to the group of automorphisms of the graded algebra A.

Let h be a splitting Cartan subalgebra of g, and R the root system of \((\mathfrak{g}, \mathfrak{h})\) . Recall that P(R) (resp. Q(R)) denotes the group of weights (resp. radical weights) of R. Put

\[
\mathrm{T} _ {\mathrm{P}} = \operatorname{Hom} (\mathrm{P} (\mathrm{R}), k ^ {*}) \quad \mathrm{T} _ {\mathrm{Q}} = \operatorname{Hom} (\mathrm{Q} (\mathrm{R}), k ^ {*}).
\]

We can consider \(g = g^{0} + \sum_{\alpha \in R} g^{\alpha}\) as a \(Q(R)\) -graded algebra. The preceding remarks define a canonical homomorphism from \(T_{Q}\) to \(\operatorname{Aut}(\mathfrak{g}, \mathfrak{h})\) , which will be denoted by f in this paragraph. In the other hand, the canonical injection from \(Q(R)\) to \(P(R)\) defines a homomorphism from \(T_{P}\) to \(T_{Q}\) , which will be denoted by q:

\[
\mathrm{T} _ {\mathrm{P}} \quad \stackrel {{q}} {{\longrightarrow}} \quad \mathrm{T} _ {\mathrm{Q}} \quad \stackrel {{f}} {{\longrightarrow}} \quad \operatorname{Aut} (\mathfrak {g}, \mathfrak {h}).
\]

If \(s \in \operatorname{Aut}(\mathfrak{g}, \mathfrak{h})\) , let \(s^*\) be the restriction of \({}^t(s|_{\mathfrak{h}})^{-1}\) to \(\mathrm{Q}(\mathrm{R})\) . Then, for all \(\varphi \in \mathrm{T}_{\mathrm{Q}}\) ,

\[
f (\varphi \circ s ^ {*}) = s ^ {- 1} \circ f (\varphi) \circ s.\tag{1}
\]

Indeed, let \(\gamma \in \mathrm{Q}(\mathrm{R})\) and \(x\in \mathfrak{g}^{\gamma}\) ; then \(sx\in \mathfrak{g}^{s^{*}\gamma}\) and

\[
f (\varphi \circ s ^ {*}) x = (\varphi \circ s ^ {*}) (\gamma). x = s ^ {- 1} (\varphi (s ^ {*} \gamma) s x) = (s ^ {- 1} \circ f (\varphi) \circ s) (x).
\]

PROPOSITION 2. The sequence of homomorphisms

\[
1 \longrightarrow \mathrm {T_ {Q}} \stackrel {f} {\longrightarrow} \mathrm{Aut} (\mathfrak {g}, \mathfrak {h}) \stackrel {\varepsilon} {\longrightarrow} \mathrm{A(R)} \longrightarrow 1
\]

is exact.

\begin{enumerate}
\def\labelenumi{\alph{enumi})}
\item
  Let \(\varphi \in \operatorname{Ker} f\) . Then \(\varphi(\alpha) = 1\) for all \(\alpha \in \mathrm{R}\) . Since \(\mathrm{R}\) generates the group \(\mathrm{Q}(\mathrm{R})\) , \(\varphi\) is the identity element of \(\mathrm{T}_{\mathrm{Q}}\) .
\item
  Let \(\varphi \in \mathrm{T}_{\mathrm{Q}}\) . The restriction of \(f(\varphi)\) to \(\mathfrak{h} = \mathfrak{g}^0\) is the identity, so
\end{enumerate}

\(\operatorname{Im} f \subset \operatorname{Ker} \varepsilon\) .

\begin{enumerate}
\def\labelenumi{\alph{enumi})}
\setcounter{enumi}{2}
\tightlist
\item
  Let \(s \in \operatorname{Ker} \varepsilon\) . Then \(s|_{\mathfrak{h}} = \mathrm{Id}_{\mathfrak{h}}\) . For all \(\alpha \in R\) , we have \(s(\mathfrak{g}^{\alpha}) = \mathfrak{g}^{\alpha}\) , and there exists a \(t_{\alpha} \in k^{*}\) such that \(sx = t_{\alpha}x\) for all \(x \in g^{\alpha}\) . Writing down the condition that \(s \in \operatorname{Aut}(\mathfrak{g})\) , we obtain the relations
\end{enumerate}

\[
t _ {\alpha} t _ {- \alpha} = 1 \qquad \text{for all } \alpha \in \mathrm{R}
\]

\[
t _ {\alpha} t _ {\beta} = t _ {\alpha + \beta} \quad \text { when } \alpha , \beta , \alpha + \beta \in \mathrm{R}.
\]

Under these conditions, there exists \(\varphi \in \mathrm{T}_{\mathrm{Q}}\) such that \(\varphi (\alpha) = t_{\alpha}\) for all \(\alpha \in \mathrm{R}\) (Chap. VI, §1, no. 6, Cor. 2 of Prop. 19). Then \(s = f(\varphi)\) . Hence, \(\operatorname{Ker} \varepsilon \subset \operatorname{Im} f\) .

\begin{enumerate}
\def\labelenumi{\alph{enumi})}
\setcounter{enumi}{3}
\tightlist
\item
  The image of \(\operatorname{Aut}(\mathfrak{g},\mathfrak{h})\) under \(\varepsilon\) contains \(\mathrm{W}(\mathrm{R})\) by §2, no. 2, Cor. of Th. 2, and contains \(\operatorname{Aut}(\mathrm{R},\mathrm{B})\) by Prop. 1. Hence this image is equal to \(\mathrm{A}(\mathrm{R})\) .
\end{enumerate}

COROLLARY 1. Let \((\mathrm{B}, (X_{\alpha})_{\alpha \in \mathrm{B}})\) be a framing of \((\mathfrak{g}, \mathfrak{h})\) . Let \(\mathrm{G}\) be the set of \(s \in \operatorname{Aut}(\mathfrak{g}, \mathfrak{h})\) that leave the framing invariant. Then \(\operatorname{Aut}(\mathfrak{g}, \mathfrak{h})\) is the semi-direct product of \(\mathrm{G}\) and \(\varepsilon^{-1}(\mathrm{W}(\mathrm{R}))\) .

Indeed, \(\mathrm{G} \cap \varepsilon^{-1}(\mathrm{W}(\mathrm{R})) = \{1\}\) by Prop. 1, and

\[
\operatorname{Aut} (\mathfrak {g}, \mathfrak {h}) = \mathrm{G}. \varepsilon^ {- 1} (\mathrm{W(R)})
\]

since \(\varepsilon\) is surjective (Prop. 2).

COROLLARY 2. The group \(\varepsilon^{-1}(\mathrm{W}(\mathrm{R}))\) operates simply-transitively on the set of framings of \((\mathfrak{g},\mathfrak{h})\) .

Indeed, \(\operatorname{Aut}(\mathfrak{g},\mathfrak{h})\) operates transitively on the set of framings of \((\mathfrak{g},\mathfrak{h})\) by §4, no. 4, Th. 2. Cor. 2 now follows from Cor. 1.

COROLLARY 3. Let B be a basis of R. The group \(\operatorname{Ker} \varepsilon = f(\mathrm{T}_{\mathrm{Q}})\) operates simply-transitively on the set of framings of \((\mathfrak{g}, \mathfrak{h})\) of the form \((\mathrm{B}, (X_{\alpha})_{\alpha \in \mathrm{B}})\) .

This follows immediately from Prop. 2.

Let \(\alpha \in \mathrm{R}\) , \(X_{\alpha} \in \mathfrak{g}^{\alpha}\) , \(X_{-\alpha} \in \mathfrak{g}^{-\alpha}\) be such that \([X_{\alpha}, X_{-\alpha}] = -H_{\alpha}\) . We have seen (§2, no. 2, Th. 2) that, for all \(t \in k^{*}\) , the restriction of the elementary automorphism

\[
\theta_ {\alpha} (t) = e ^ {\mathrm{ad} t X _ {\alpha}} e ^ {\mathrm{ad} t ^ {- 1} X _ {- \alpha}} e ^ {\mathrm{ad} t X _ {\alpha}}
\]

to \(\mathfrak{h}\) is the transpose of \(s_{\alpha}\) ; so \(\varepsilon(\theta_{\alpha}(t)) = s_{\alpha}\) and consequently \(\theta_{\alpha}(t)\theta_{\alpha}(-1) \in \operatorname{Ker} \varepsilon\) .

Lemma 1. Let \(\alpha \in \mathrm{R}\) and \(t\in k^{*}\) . Let \(\varphi\) be the homomorphism \(\lambda \mapsto t^{\lambda (H_{\alpha})}\) from \(\mathrm{Q(R)}\) to \(k^*\) . Then \(f(\varphi) = \theta_{\alpha}(t)\theta_{\alpha}(-1)\) .

Let \(\rho\) be the representation of \(\mathfrak{sl}(2,k)\) on \(\mathfrak{g}\) associated to \(X_{\alpha}\) . Let \(\pi\) be the representation of \(\mathbf{SL}(2,k)\) compatible with \(\rho\) . Introduce the notations \(\theta(t), h(t)\) of §1, no. 5. Since \(\rho(H) = \operatorname{ad} H_{\alpha}\) , the elements of \(\mathfrak{g}^{\lambda}\) are of weight \(\lambda(H_{\alpha})\) for \(\rho\) . By §2, no. 2, \(\theta_{\alpha}(t)\theta_{\alpha}(-1) = \pi(\theta(t)\theta(-1)) = \pi(h(t))\) . Hence the restriction of \(\theta_{\alpha}(t)\theta_{\alpha}(-1)\) to \(\mathfrak{g}^{\lambda}\) is the homothety of ratio \(t^{\lambda(H_{\alpha})}\) (§1, no. 5, Prop. 6), hence the lemma.

PROPOSITION 3. The image of the composite homomorphism

\[
\mathrm{T} _ {\mathrm{P}} \quad \xrightarrow {q} \quad \mathrm{T} _ {\mathrm{Q}} \quad \xrightarrow {f} \quad \operatorname{Aut} (\mathfrak {g}, \mathfrak {h})
\]

is contained in \(\mathrm{Aut}_e(\mathfrak{g})\) .

Let B be a basis of R. Then \((H_{\alpha})_{\alpha\in\mathrm{B}}\) is a basis of \(R^{\vee}\) , and the dual basis of \((H_{\alpha})_{\alpha\in\mathrm{B}}\) in \(h^{*}\) is a basis of the group P(R). Hence the group \(T_{P}\) is generated by the homomorphisms \(\lambda\mapsto t^{\lambda(H_{\alpha})}\;(t\in k^{*},\alpha\in\mathrm{B})\) . If \(\varphi\) is the restriction of such a homomorphism to Q(R), Lemma 1 proves that \(f(\varphi)\in\operatorname{Aut}_{e}(\mathfrak{g})\) , hence the proposition.

Let \(\bar{k}\) be an algebraic closure of \(k\) . The map which associates to any automorphism \(s\) of \(\mathfrak{g}\) the automorphism \(s \otimes 1\) of \(\mathfrak{g} \otimes_k \bar{k}\) is an injective homomorphism from \(\mathrm{Aut}(\mathfrak{g})\) to \(\mathrm{Aut}(\mathfrak{g} \otimes_k \bar{k})\) . We denote by \(\mathrm{Aut}_0(\mathfrak{g})\) the normal subgroup of \(\mathrm{Aut}(\mathfrak{g})\) which is the inverse image of \(\mathrm{Aut}_e(\mathfrak{g} \otimes_k \bar{k})\) under this homomorphism; this is the set of automorphisms of \(\mathfrak{g}\) that become elementary on extending the base field from \(k\) to \(\bar{k}\) . It is clear that \(\mathrm{Aut}_e(\mathfrak{g})\) is independent of the choice of \(\bar{k}\) , and that \(\mathrm{Aut}_e(\mathfrak{g}) \subset \mathrm{Aut}_0(\mathfrak{g})\) . The groups \(\mathrm{Aut}_0(\mathfrak{g})\) and \(\mathrm{Aut}_e(\mathfrak{g})\) can be distinct (Chap. VII, §13, no. 1). If \(\mathfrak{h}\) is a Cartan subalgebra of \(\mathfrak{g}\) , put

\[
\operatorname{Aut} _ {e} (\mathfrak {g}, \mathfrak {h}) = \operatorname{Aut} _ {e} (\mathfrak {g}) \cap \operatorname{Aut} (\mathfrak {g}, \mathfrak {h}), \quad \operatorname{Aut} _ {0} (\mathfrak {g}, \mathfrak {h}) = \operatorname{Aut} _ {0} (\mathfrak {g}) \cap \operatorname{Aut} (\mathfrak {g}, \mathfrak {h}).
\]

Lemma 2. Let \(\mathfrak{h}\) be a splitting Cartan subalgebra of \(\mathfrak{g}\) , and \(s \in \operatorname{Aut}_0(\mathfrak{g}, \mathfrak{h})\) . Assume that the restriction of \(s\) to \(\sum_{\alpha \in \mathrm{R}} \mathfrak{g}^\alpha\) does not have 1 as an eigenvalue. Then \(\varepsilon(s) = 1\) .

By extension of \(k\) , we are reduced to the case where \(s \in \mathrm{Aut}_e(\mathfrak{g}, \mathfrak{h})\) . The dimension of the nilspace of \(s - 1\) is at least \(\dim \mathfrak{h}\) (Chap. VII, §4, no. 4, Prop. 9). Hence \((s - 1)|_{\mathfrak{h}}\) is nilpotent. Since \(s|_{\mathfrak{h}} \in \mathrm{A}(\mathrm{R}^\vee)\) , \(s|_{\mathfrak{h}}\) is of finite order, and hence semi-simple (Chap. V, Appendix, Prop. 2). Consequently, \((s - 1)|_{\mathfrak{h}} = 0\) , which proves that \(\varepsilon(s) = 1\) .

Lemma 3. (i) Let \(m = (\mathrm{P}(\mathrm{R}) : \mathrm{Q}(\mathrm{R}))\) . If \(\varphi\) is the \(m\) th power of an element of \(\mathrm{T}_{\mathrm{Q}}\) , then \(\varphi \in q(\mathrm{T}_{\mathrm{P}})\) .

\begin{enumerate}
\def\labelenumi{(\roman{enumi})}
\setcounter{enumi}{1}
\tightlist
\item
  If \(k\) is algebraically closed, \(q(\mathrm{T}_{\mathrm{P}}) = \mathrm{T}_{\mathrm{Q}}\) .
\end{enumerate}

There exist a basis \((\lambda_{1},\ldots,\lambda_{l})\) of P(R) and integers \(n_{1}\geq1,\ldots,n_{l}\geq1\) such that \((n_{1}\lambda_{1},\ldots,n_{l}\lambda_{l})\) is a basis of Q(R). We have \(m=n_{1}\ldots n_{l}\) . Let \(\psi\in T_{Q}\) and put \(\psi(n_{1}\lambda_{1})=t_{1},\ldots,\psi(n_{l}\lambda_{l})=t_{l}\) . For \(i=1,\ldots,l\) , put \(m_{i}=\) \(\prod_{j \neq i} n_j\) . Let \(\chi\) be the element of \(T_P\) such that \(\chi(\lambda_1) = t_1^{m_1}, \ldots, \chi(\lambda_l) = t_l^{m_l}\) . Then

\[
\chi (n _ {i} \lambda_ {i}) = t _ {i} ^ {m _ {i} n _ {i}} = t _ {i} ^ {m} = (\psi^ {m}) (n _ {i} \lambda_ {i})
\]

so \(\chi|_{\mathrm{Q}(\mathrm{R})}=\psi^{m}\) . This proves (i). If k is algebraically closed, every element of \(k^{*}\) is the mth power of an element of \(k^{*}\) , so every element of \(T_{Q}\) is the mth power of an element of \(T_{Q}\) ; hence, (ii) follows from (i).

PROPOSITION 4. We have \(f(\mathrm{T}_{\mathrm{Q}}) \subset \mathrm{Aut}_0(\mathfrak{g}, \mathfrak{h})\) and \(\varepsilon^{-1}(\mathrm{W}(\mathrm{R})) = \mathrm{Aut}_0(\mathfrak{g}, \mathfrak{h})\) .

\begin{enumerate}
\def\labelenumi{\alph{enumi})}
\tightlist
\item
  Let \(\varphi \in \mathrm{T}_{\mathrm{Q}}\) and let \(\bar{k}\) be an algebraic closure of \(k\) . By Lemma 3, \(\varphi\) extends to an element of \(\operatorname{Hom}(\mathrm{P}(\mathrm{R}), k^*)\) . By Prop. 3,
\end{enumerate}

\[
f (\varphi) \otimes 1 \in \operatorname{Aut} _ {e} (\mathfrak {g} \otimes_ {k} \bar {k}, \mathfrak {h} \otimes_ {k} \bar {k}).
\]

Hence \(f(\varphi) \in \mathrm{Aut}_0(\mathfrak{g},\mathfrak{h})\) , and \(\operatorname{Ker} \varepsilon \subset \mathrm{Aut}_0(\mathfrak{g},\mathfrak{h})\) .

\begin{enumerate}
\def\labelenumi{\alph{enumi})}
\setcounter{enumi}{1}
\item
  The image of \(\mathrm{Aut}_e(\mathfrak{g},\mathfrak{h})\) under \(\varepsilon\) contains \(\mathrm{W(R)}\) (§2, no. 2, Cor. of Th. 2). In view of a), we see that \(\varepsilon^{-1}(\mathrm{W(R)})\subset \mathrm{Aut}_0(\mathfrak{g},\mathfrak{h})\) .
\item
  It remains to prove that \(\mathrm{Aut}_0(\mathfrak{g},\mathfrak{h})\subset \varepsilon^{-1}(\mathrm{W}(\mathrm{R}))\) . In view of b), it suffices to prove that \(\varepsilon (\mathrm{Aut}_0(\mathfrak{g},\mathfrak{h}))\cap \mathrm{Aut}(\mathrm{R},\mathrm{B})\) , where B denotes a basis of R, reduces to \{1\}.
\end{enumerate}

Let \(s \in \mathrm{Aut}_0(\mathfrak{g}, \mathfrak{h})\) be such that \(\varepsilon(s) \in \mathrm{Aut}(\mathrm{R}, \mathrm{B})\) . The subgroup of \(\mathrm{A}(\mathrm{R})\) generated by \(\varepsilon(s)\) has a finite number of orbits on \(\mathrm{R}\) . Let \(\mathrm{U}\) be such an orbit, of cardinal \(r\) , and \(\mathfrak{g}^{\mathrm{U}} = \sum_{\beta \in \mathrm{U}} \mathfrak{g}^{\beta}\) . Let \(\beta_1 \in \mathrm{U}\) , and put \(\beta_i = \varepsilon(s)^{i-1}\beta_1\) for \(1 \leq i \leq r\) , so that \(\mathrm{U} = \{\beta_1, \ldots, \beta_r\}\) . Let \(X_{\beta_1}\) be a non-zero element of \(\mathfrak{g}^{\beta_1}\) , and put \(X_{\beta_i} = s^{i-1}X_{\beta_1}\) for \(1 \leq i \leq r\) . There exists \(c_{\mathrm{U}} \in k^*\) such that \(s^r X_{\beta_1} = c_{\mathrm{U}}X_{\beta_1}\) , hence \(s^r X_{\beta_i} = c_{\mathrm{U}}X_{\beta_i}\) for all \(i\) , and consequently \(s^r |_{\mathfrak{g}^{\mathrm{U}}} = c_{\mathrm{U}}.1\) . Let \(\varphi \in T_{\mathrm{Q}}\) , and \(s' = s \circ f(\varphi)\) , which by \(a)\) is an element of \(\mathrm{Aut}_0(\mathfrak{g}, \mathfrak{h})\) . We have \(s'^r |_{\mathfrak{g}^{\mathrm{U}}} = c_{\mathrm{U}}'.1\) , where

\[
c _ {\mathrm{U}} ^ {\prime} = c _ {\mathrm{U}} \prod_ {i = 1} ^ {r} \varphi (\beta_ {i}) = c _ {\mathrm{U}} \varphi \left(\sum_ {i = 1} ^ {r} \beta_ {i}\right).
\]

Put \(\mathrm{B} = \{\alpha_1, \ldots, \alpha_l\}\) and \(\sum_{i=1}^{r} \beta_i = \sum_{j=1}^{l} m_j^{\mathrm{U}} \alpha_j\) . Since \(\varepsilon(s) \in \mathrm{Aut}(\mathrm{R}, \mathrm{B})\) , the \(m_j^{\mathrm{U}}\) are integers of the same sign and not all zero. We have

\[
c _ {\mathrm{U}} ^ {\prime} = c _ {\mathrm{U}} \prod_ {j = 1} ^ {l} \varphi (\alpha_ {j}) ^ {m _ {j} ^ {\mathrm{U}}}.
\]

Now \(\varphi\) can be chosen so that \(c_{\mathrm{U}}^{\prime} \neq 1\) for every orbit U; indeed, this reduces to choosing elements \(\varphi(\alpha_1) = t_1, \ldots, \varphi(\alpha_l) = t_l\) of \(k^*\) which are not annihilated by a finite number of polynomials in \(t_1, \ldots, t_l\) , not identically zero. For such a choice of \(\varphi\) , \(\varepsilon(s') = 1\) by Lemma 2, so

\[
\varepsilon (s) = \varepsilon (s ^ {\prime}) \varepsilon (f (\varphi)) ^ {- 1} = 1.
\]

COROLLARY. Let B be a basis of R. The group \(\operatorname{Aut}(\mathfrak{g},\mathfrak{h})\) is isomorphic to the semi-direct product of the groups \(\operatorname{Aut}(\mathrm{R},\mathrm{B})\) and \(\operatorname{Aut}_{0}(\mathfrak{g},\mathfrak{h})\) .

This follows from Prop. 1, Cor. 1 of Prop. 2, and Prop. 4.

Remark. Let \(\varepsilon^{\prime},\varepsilon^{\prime\prime}\) be the restrictions of \(\varepsilon\) to \(\operatorname{Aut}_{0}(\mathfrak{g},\mathfrak{h}),\operatorname{Aut}_{e}(\mathfrak{g},\mathfrak{h})\) . Let \(f^{\prime}\) be the homomorphism from \(T_{P}\) to \(\operatorname{Aut}_{e}(\mathfrak{g},\mathfrak{h})\) induced by f via the canonical injection from \(Q(R)\) to \(P(R)\) . In the preceding we have established the following commutative diagram:

\[
\begin{array}{c c c c c c c c c} 1 & \longrightarrow & \mathrm{T} _ {\mathrm{Q}} & \stackrel {{f}} {{\longrightarrow}} & \mathrm{Aut} (\mathfrak {g}, \mathfrak {h}) & \stackrel {{\varepsilon}} {{\longrightarrow}} & \mathrm{A} (\mathrm{R}) & \longrightarrow & 1 \\ & & \uparrow & & \uparrow & & \uparrow & & \\ 1 & \longrightarrow & \mathrm{T} _ {\mathrm{Q}} & \stackrel {{f}} {{\longrightarrow}} & \mathrm{Aut} _ {0} (\mathfrak {g}, \mathfrak {h}) & \stackrel {{\varepsilon^ {\prime}}} {{\longrightarrow}} & \mathrm{W} (\mathrm{R}) & \longrightarrow & 1 \\ & & \uparrow_ {q} & & \uparrow & & \uparrow & & \\ & & \mathrm{T} _ {\mathrm{P}} & \stackrel {{f ^ {\prime}}} {{\longrightarrow}} & \mathrm{Aut} _ {e} (\mathfrak {g}, \mathfrak {h}) & \stackrel {{\varepsilon^ {\prime \prime}}} {{\longrightarrow}} & \mathrm{W} (\mathrm{R}) & \longrightarrow & 1 \end{array}
\]

in which the vertical arrows other than q denote the canonical injections. We have seen (Prop. 2 and 4) that the first two rows are exact. In the third row, the homomorphism \(\varepsilon''\) is surjective (§2, no. 2, Cor. of Th. 2); it can be shown that its kernel is \(f'(\mathrm{T}_{\mathrm{P}})\) (§7, Exerc. 26 d)).

\subsection{AUTOMORPHISMS OF A SPLITTABLE SEMI-SIMPLE LIE ALGEBRA}

PROPOSITION 5. Assume that g is splittable. The group \(\operatorname{Aut}_{0}(\mathfrak{g})\) operates simply-transitively on the set of framings of g.

Let \(e_1 = (\mathfrak{g}, \mathfrak{h}_1, \mathrm{B}_1, (X_\alpha^1)_{\alpha \in \mathrm{B}_1})\) , \(e_2 = (\mathfrak{g}, \mathfrak{h}_2, \mathrm{B}_2, (X_\alpha^2)_{\alpha \in \mathrm{B}_2})\) be two framings of \(\mathfrak{g}\) . There exists at least one element of \(\operatorname{Aut}_0(\mathfrak{g})\) that transforms \(e_1\) into \(e_2\) (Prop. 1 and Prop. 4). Let \(\bar{k}\) be an algebraic closure of \(k\) . There exists an element of \(\operatorname{Aut}_e(\mathfrak{g} \otimes_k \bar{k})\) that transforms \(\mathfrak{h}_1 \otimes_k \bar{k}\) into \(\mathfrak{h}_2 \otimes_k \bar{k}\) (Chap. VII, §3, no. 2, Th. 1). Hence, by Prop. 4 and Cor. 2 of Prop. 2, there exists an element \(\varphi\) of \(\operatorname{Aut}_e(\mathfrak{g} \otimes_k \bar{k})\) that transforms the framing \((\mathfrak{g} \otimes_k \bar{k}, \mathfrak{h}_1 \otimes_k \bar{k}, \mathrm{B}_1, (X_\alpha^1)_{\alpha \in \mathrm{B}_1})\) of \(\mathfrak{g} \otimes_k \bar{k}\) into the framing \((\mathfrak{g} \otimes_k \bar{k}, \mathfrak{h}_2 \otimes_k \bar{k}, \mathrm{B}_2, (X_\alpha^2)_{\alpha \in \mathrm{B}_2})\) . Since \(\mathfrak{h}_1\) and the \(X_\alpha^1\) (resp. \(\mathfrak{h}_2\) and the \(X_\alpha^2\) ) generate \(\mathfrak{g}_1\) (resp. \(\mathfrak{g}_2\) ), we have \(\varphi(\mathfrak{g}_1) = \mathfrak{g}_2\) , so \(\varphi\) is of the form \(\psi \otimes 1\) where \(\psi \in \operatorname{Aut}_0(\mathfrak{g})\) , and \(\psi\) transforms \(e_1\) into \(e_2\) .

COROLLARY 1. Let \((\mathfrak{g},\mathfrak{h},\mathrm{B},(X_{\alpha})_{\alpha\in\mathrm{B}})\) be a framing of g, and G the group (isomorphic to \(\operatorname{Aut}(\mathrm{R},\mathrm{B})\) ) of automorphisms of g that leave this framing invariant. Then \(\operatorname{Aut}(\mathfrak{g})\) is the semi-direct product of G and \(\operatorname{Aut}_{0}(\mathfrak{g})\) .

Indeed, every element of \(\operatorname{Aut}(\mathfrak{g})\) transforms \((\mathfrak{g},\mathfrak{h},\mathrm{B},(X_{\alpha})_{\alpha\in\mathrm{B}})\) into a framing of g. By Prop. 5, every coset of \(\operatorname{Aut}(\mathfrak{g})\) modulo \(\operatorname{Aut}_{0}(\mathfrak{g})\) meets G in exactly one point. Q.E.D.

It follows from Cor. 1 that the group \(\operatorname{Aut}(\mathfrak{g})/\operatorname{Aut}_{0}(\mathfrak{g})\) can be identified with \(\operatorname{Aut}(\mathrm{R},\mathrm{B})\) , and is isomorphic to the group of automorphisms of the Dynkin graph of R.

COROLLARY 2. \(\operatorname{Aut}(\mathfrak{g}) = \operatorname{Aut}_0(\mathfrak{g})\) when \(\mathfrak{g}\) is a splittable simple Lie algebra of type \(\mathrm{A}_1\) , \(\mathrm{B}_n\) ( \(n \geq 2\) ), \(\mathrm{C}_n\) ( \(n \geq 2\) ), \(\mathrm{E}_7\) , \(\mathrm{E}_8\) , \(\mathrm{F}_4\) , \(\mathrm{G}_2\) . The quotient \(\operatorname{Aut}(\mathfrak{g}) / \operatorname{Aut}_0(\mathfrak{g})\) is of order 2 when \(\mathfrak{g}\) is of type \(\mathrm{A}_n\) ( \(n \geq 2\) ), \(\mathrm{D}_n\) ( \(n \geq 5\) ), \(\mathrm{E}_6\) ; it is isomorphic to \(\mathfrak{S}_3\) when \(\mathfrak{g}\) is of type \(\mathrm{D}_4\) .

This follows from Cor. 1 and Chap. VI, Plates I to IX.

Remarks. 1) Let \(e_1 = (\mathfrak{g}, \mathfrak{h}_1, \mathrm{B}_1, (X_\alpha^1)_{\alpha \in \mathrm{B}_1})\) , \(e_2 = (\mathfrak{g}, \mathfrak{h}_2, \mathrm{B}_2, (X_\alpha^2)_{\alpha \in \mathrm{B}_2})\) , \(e_2' = (\mathfrak{g}, \mathfrak{h}_2, \mathrm{B}_2, (Y_\alpha^2)_{\alpha \in \mathrm{B}_2})\) be framings of \(\mathfrak{g}\) , and \(s\) (resp. \(s'\) ) an element of \(\operatorname{Aut}_0(\mathfrak{g})\) that transforms \(e_1\) to \(e_2\) (resp. \(e_2'\) ). Then \(s|_{\mathfrak{h}_1} = s'|_{\mathfrak{h}_1}\) . Indeed, \(s'^{-1}s \in \operatorname{Aut}_0(\mathfrak{g}, \mathfrak{h}_1)\) and \(s'^{-1}s(\mathrm{B}_1) = \mathrm{B}_1\) , so \(\varepsilon(s'^{-1}s) = 1\) .

\begin{enumerate}
\def\labelenumi{\arabic{enumi})}
\setcounter{enumi}{1}
\tightlist
\item
  Let X be the set of pairs \((\mathfrak{h},\mathrm{B})\) where h is a splitting Cartan subalgebra of g and B a basis of the root system of \((\mathfrak{g},\mathfrak{h})\) . If \(x=(\mathfrak{h},\mathrm{B})\) and \(x'=(\mathfrak{h}',\mathrm{B}')\) are two elements of X, there exists \(s\in\operatorname{Aut}_{0}(\mathfrak{g})\) that transforms x into \(x'\) (Prop. 5), and the restriction \(s_{x',x}\) of s to h does not depend on the choice of s (Remark 1). In particular, \(s_{x'',x'}\circ s_{x',x}=s_{x'',x}\) if \(x,x',x''\in X\) , and \(s_{x,x}=1\) . The set of families \((h_{x})_{x\in\mathrm{X}}\) satisfying the conditions
\end{enumerate}

\[
a) h _ {x} \in \mathfrak {h} \text {   if   } x = (\mathfrak {h}, B)
\]

\[
b) s _ {x ^ {\prime}, x} (h _ {x}) = h _ {x ^ {\prime}} \text { if } x, x ^ {\prime} \in X
\]

is in a natural way a vector space \(\mathfrak{h}(\mathfrak{g})\) which we sometimes call the canonical Cartan subalgebra of g. For \(x = (\mathfrak{h}, \mathrm{B})\) and \(x' = (\mathfrak{h}', \mathrm{B}')\) , \(s_{x', x}\) takes B to \(B'\) , and hence the root system of \((\mathfrak{g}, \mathfrak{h})\) to that of \((\mathfrak{g}, \mathfrak{h}')\) ; it follows that the dual \(\mathfrak{h}(\mathfrak{g})^*\) of \(\mathfrak{h}(\mathfrak{g})\) is naturally equipped with a root system \(\mathrm{R}(\mathfrak{g})\) and with a basis \(\mathrm{B}(\mathfrak{g})\) of \(\mathrm{R}(\mathfrak{g})\) . We sometimes say that \(\mathrm{R}(\mathfrak{g})\) is the canonical root system of g and that \(\mathrm{B}(\mathfrak{g})\) is its canonical basis. The group \(\mathrm{Aut}(\mathfrak{g})\) operates on \(\mathfrak{h}(\mathfrak{g})\) leaving \(\mathrm{R}(\mathfrak{g})\) and \(\mathrm{B}(\mathfrak{g})\) stable; the elements of \(\mathrm{Aut}(\mathfrak{g})\) that operate trivially on \(\mathfrak{h}(\mathfrak{g})\) are those of \(\mathrm{Aut}_0(\mathfrak{g})\) .

PROPOSITION 6. Let \(\mathfrak{h}\) be a splitting Cartan subalgebra of \(\mathfrak{g}\) . We have, with the notations in no. 1, \(\mathrm{Aut}_0(\mathfrak{g}) = \mathrm{Aut}_e(\mathfrak{g}).\mathrm{Ker}\varepsilon = \mathrm{Aut}_e(\mathfrak{g}).f(\mathrm{T}_Q)\) .

By §3, no. 3, Cor. of Prop. 10, \(\mathrm{Aut}_0(\mathfrak{g}) = \mathrm{Aut}_e(\mathfrak{g}).\mathrm{Aut}_0(\mathfrak{g},\mathfrak{h})\) . On the other hand, \(\varepsilon (\mathrm{Aut}_e(\mathfrak{g},\mathfrak{h}))\supset \mathrm{W}(\mathrm{R})\) by §2, no. 2, Cor. of Th. 2, so \(\mathrm{Aut}_0(\mathfrak{g},\mathfrak{h}) = \mathrm{Aut}_e(\mathfrak{g},\mathfrak{h}).\mathrm{Ker}\varepsilon\) .

Remark 3. Prop. 6 shows that the canonical homomorphism

\[
\iota : \mathrm{T} _ {\mathrm{Q}} / \operatorname{Im} (\mathrm{T} _ {\mathrm{P}}) \to \operatorname{Aut} _ {0} (\mathfrak {g}) / \operatorname{Aut} _ {e} (\mathfrak {g}),
\]

induced by the diagram in no. 2, is surjective. In particular, \(\mathrm{Aut}_e(\mathfrak{g})\) contains the derived group of \(\mathrm{Aut}_0(\mathfrak{g})\) ; we shall see (§11, no. 2, Prop. 3) that they are actually equal. Moreover, it can be shown that \(\iota\) is injective, in other words that

\[
f (\mathrm{T} _ {\mathrm{Q}}) \cap \operatorname{Aut} _ {e} (\mathfrak {g}) = f ^ {\prime} (\mathrm{T} _ {\mathrm{P}}),
\]

(cf.~§7, Exerc. 26 d)).

PROPOSITION 7. Let g be a splittable semi-simple Lie algebra, b a Borel subalgebra of g, and \(p_{1}\) and \(p_{2}\) distinct parabolic subalgebras of g containing b. Then \(p_{1}\) and \(p_{2}\) are not conjugate under \(\operatorname{Aut}_{0}(\mathfrak{g})\) .

We can assume that \(k\) is algebraically closed. Let \(s \in \mathrm{Aut}_0(\mathfrak{g})\) be such that \(s(\mathfrak{p}_1) = \mathfrak{p}_2\) . Let \(\mathfrak{h}\) be a Cartan subalgebra of \(\mathfrak{g}\) contained in \(\mathfrak{b} \cap s(\mathfrak{b})\) (§3, no. 3, Prop. 10). Since \(\mathfrak{h}\) and \(s(\mathfrak{h})\) are Cartan subalgebras of \(s(\mathfrak{b})\) , there exists \(u \in [\mathfrak{b}, \mathfrak{b}]\) such that \(e^{\mathrm{ad}u}(\mathfrak{h}) = s(\mathfrak{h})\) (Chap. VII, §3, no. 4, Th. 3). Replacing \(s\) by \(e^{-\mathrm{ad}u}s\) , we are reduced to the case in which \(s(\mathfrak{h}) = \mathfrak{h}\) , and \(s\) then induces on \(\mathfrak{h}\) an element \(\sigma\) of the Weyl group W of \((\mathfrak{g}, \mathfrak{h})\) (Prop. 4). Let C be the Weyl chamber corresponding to \(\mathfrak{b}\) . Then \(\mathfrak{p}_1\) and \(\mathfrak{p}_2\) correspond to facets \(F_1\) and \(F_2\) of \(\mathfrak{h}_{\mathbf{R}}\) contained in the closure of C. We have \(\sigma(F_1) = F_2\) . Since \(\sigma \in W\) , this implies that \(F_1 = F_2\) (Chap. V, §3, no. 3, Th. 2) so \(\mathfrak{p}_1 = \mathfrak{p}_2\) .

Remark 4. Let g be a splittable semi-simple Lie algebra, P the set of parabolic subalgebras of g, a set on which \(\operatorname{Aut}_{0}(\mathfrak{g})\) operates. Retain the notations of Remark 2. Let \(\Sigma\) be a subset of \(\mathrm{B}(\mathfrak{g})\) . Giving \(\Sigma\) is equivalent to giving, for every \(x = (\mathfrak{h}, \mathrm{B}) \in \mathrm{X}\) , a subset \(\Sigma_{x}\) of B, such that \(s_{x',x}\) takes \(\Sigma_{x}\) to \(\Sigma_{x'}\) for any \(x, x' \in X\) . Let \(p_{x}\) be the parabolic subalgebra of g corresponding to \(\Sigma_{x}\) (§3, no. 4, Remark). The orbit of \(p_{x}\) under \(\operatorname{Aut}_{0}(\mathfrak{g})\) is the set of \(p_{x'}\) for \(x' \in X\) . This defines a map from \(\mathfrak{P}(\mathrm{B}(\mathfrak{g}))\) to \(\mathcal{P}/\operatorname{Aut}_{0}(\mathfrak{g})\) . This map is surjective by the Remark of §3, no. 4, and injective by Prop. 7.

\subsection{ZARISKI TOPOLOGY ON Aut(g)}

PROPOSITION 8. Let V be the set of endomorphisms of the vector space g. Then \(\operatorname{Aut}(\mathfrak{g})\) is closed in V for the Zariski topology (Chap. VII, App. I).

Let \(K\) be the Killing form of \(\mathfrak{g}\) . If \(s \in \operatorname{Aut}(\mathfrak{g})\) ,

\[
  [ s x, s y ] = [ x, y ]\tag{2}
\]

\[
\mathrm{K} (s x, s y) = \mathrm{K} (x, y)\tag{3}
\]

for all \(x, y \in g\) . Conversely, let s be an element of V satisfying (2) and (3) for all \(x, y \in g\) . Then \(\operatorname{Ker}(s) = 0\) , so s is bijective and \(s \in \operatorname{Aut}(\mathfrak{g})\) . But, for all \(x, y \in g\) , the maps \(s \mapsto [sx, sy]\) and \(s \mapsto \operatorname{K}(sx, sy)\) from V to g and k are polynomial.

PROPOSITION 9. Let \(\mathfrak{h}\) be a splitting Cartan subalgebra of \(\mathfrak{g}\) .

\begin{enumerate}
\def\labelenumi{(\roman{enumi})}
\item
  The group \(f(\mathrm{T}_{\mathrm{Q}})\) is closed in \(\operatorname{Aut}(\mathfrak{g})\) in the Zariski topology.
\item
  The group \(f(q(\mathrm{T_P}))\) is dense in \(f(\mathrm{T_Q})\) in the Zariski topology.
\end{enumerate}

Assertion (i) follows from the equality \(f(\mathrm{T}_{\mathrm{Q}}) = \operatorname{Aut}(\mathfrak{g},\mathfrak{h})\cap \operatorname{Ker}\varepsilon\) (Prop. 2). Put \(m = (\mathrm{P}(\mathrm{R}):\mathrm{Q}(\mathrm{R}))\) . Let \(\mathrm{F}\) be a polynomial function on \(\mathrm{V}\) ; we assume that \(\mathrm{F}\) vanishes on the \(m\) th power of every element of \(f(\mathrm{T}_{\mathrm{Q}})\) , and show that \(\mathrm{F}|f(\mathrm{T}_{\mathrm{Q}}) = 0\) ; in view of Lemma 3, this will prove (ii).

The set \(V'\) of elements of V inducing the identity on h and leaving each \(g^{\alpha}\) stable can be identified with \(k^{R}\) . Let \(F'\) be the restriction of F to \(V' = k^{R}\) ; this is a polynomial function. We have \(f(\mathrm{T}_{\mathrm{Q}}) \subset \mathrm{V}'\) . Let \(\mathrm{B} = (\alpha_{1}, \ldots, \alpha_{l})\) be a basis of R. For all \(t = (t_{1}, \ldots, t_{l}) \in k^{*B}\) , let \(\varphi(t)\) be the homomorphism from \(\mathrm{Q}(\mathrm{R})\) to the group \(k^{*}\) that extends t. Then \(\mathrm{F}'(f(\varphi(t)))\) can be written as a finite sum

\[
\sum_ {n _ {1}, \dots , n _ {l} \in \mathbf {Z}} c _ {n _ {1}, \dots , n _ {l}} t _ {1} ^ {n _ {1}} \dots t _ {l} ^ {n _ {l}} = \mathrm{H} (t _ {1}, \dots , t _ {l}).
\]

By assumption,

\[
0 = \mathrm{H} (t _ {1} ^ {m}, \dots , t _ {l} ^ {m}) = \sum_ {n _ {1}, \dots , n _ {l} \in \mathbf {Z}} c _ {n _ {1}, \dots , n _ {l}} t _ {1} ^ {m n _ {1}} \dots t _ {l} ^ {m n _ {l}}
\]

for all \(t_{1},\ldots,t_{l}\in k^{*}\) . The \(c_{n_{1},\ldots,n_{l}}\) are thus the coefficients of a polynomial in l variables which vanishes on \(k^{*l}\) ; hence they are all zero.

PROPOSITION 10. Assume that \(\mathfrak{g}\) is splittable.

\begin{enumerate}
\def\labelenumi{(\roman{enumi})}
\item
  The group \(\operatorname{Aut}_{e}(\mathfrak{g})\) is dense in \(\operatorname{Aut}_{0}(\mathfrak{g})\) in the Zariski topology.
\item
  The groups \(\mathrm{Aut}_e(\mathfrak{g})\) and \(\mathrm{Aut}_0(\mathfrak{g})\) are connected in the Zariski topology.
\end{enumerate}

By Prop. 3, \(f(q(\mathrm{T}_{\mathrm{P}})) \subset \mathrm{Aut}_e(\mathfrak{g})\) . For all \(s \in \mathrm{Aut}_e(\mathfrak{g})\) , the closure of \(s.f(q(\mathrm{T}_{\mathrm{P}}))\) in the Zariski topology contains \(s.f(\mathrm{T}_{\mathrm{Q}})\) by Prop. 9. Hence the closure of \(\mathrm{Aut}_e(\mathfrak{g})\) contains \(\mathrm{Aut}_e(\mathfrak{g}).f(\mathrm{T}_{\mathrm{Q}}) = \mathrm{Aut}_0(\mathfrak{g})\) (Prop. 6). This proves (i).

Let \(\mathrm{Aut}_e(\mathfrak{g}) = \Omega \cup \Omega'\) be a partition of \(\mathrm{Aut}_e(\mathfrak{g})\) formed by relatively open subsets in the Zariski topology, and with \(\Omega \neq \emptyset\) . If \(\omega \in \Omega\) and if \(x\) is a nilpotent element of \(\mathfrak{g}\) , the map \(\tau : t \mapsto \omega \exp(t \operatorname{ad} x)\) from \(k\) to \(\mathrm{Aut}_e(\mathfrak{g})\) is polynomial, hence continuous in the Zariski topology; consequently, \(\tau(k)\) is connected; since \(\omega \in \tau(k)\) , we have \(\tau(k) \subset \Omega\) . Thus, \(\Omega. (\exp \operatorname{ad} kx) \subset \Omega\) , so \(\Omega. \mathrm{Aut}_e(\mathfrak{g}) \subset \Omega\) and \(\Omega = \mathrm{Aut}_e(\mathfrak{g})\) . This proves that \(\mathrm{Aut}_e(\mathfrak{g})\) is connected. It follows, by (i), that \(\mathrm{Aut}_0(\mathfrak{g})\) is connected. Q.E.D.

We shall see (§8, no. 4, Cor. of Prop. 6) that \(\mathrm{Aut}_0(\mathfrak{g})\) is closed in V in the Zariski topology, and that it is the connected component of the identity element of \(\mathrm{Aut}(\mathfrak{g})\) . On the other hand, \(\mathrm{Aut}_e(\mathfrak{g})\) is not in general closed in the Zariski topology.

* Assume that \((\mathfrak{g},\mathfrak{h})\) is split. The group \(\mathrm{Aut}_0(\mathfrak{g})\) is the group \(G(k)\) of \(k\) -points of a connected semi-simple algebraic group \(G\) with trivial centre (adjoint group). The group \(f(T_{Q})\) is equal to \(H(k)\) , where \(H\) is the Cartan subgroup of \(G\) with Lie algebra \(\mathfrak{h}\) . The inverse image \(\tilde{H}\) of \(H\) in the universal covering \(\tilde{G}\) of \(G\) (in the algebraic sense) has \(T_{P}\) as its group of \(k\) -points. The image of \(\tilde{G}(k)\) in \(G(k) = \mathrm{Aut}_0(\mathfrak{g})\) is the group \(\mathrm{Aut}_e(\mathfrak{g})\) .*

\subsection{LIE GROUP CASE}

PROPOSITION 11. Assume that k is R, C or a non-discrete complete ultrametric field. Let h be a splitting Cartan subalgebra of g.

\begin{enumerate}
\def\labelenumi{(\roman{enumi})}
\item
  \(\operatorname{Aut}(\mathfrak{g},\mathfrak{h})\) is a Lie subgroup of \(\operatorname{Aut}(\mathfrak{g})\) with Lie algebra ad h.
\item
  \(f(\mathrm{T}_{\mathrm{Q}})\) and \((q\circ f)(\mathrm{T}_{\mathrm{P}})\) are open subgroups of \(\operatorname {Aut}(\mathfrak{g},\mathfrak{h})\)
\item
  \(\mathrm{Aut}_e(\mathfrak{g})\) is an open subgroup of \(\mathrm{Aut}(\mathfrak{g})\) .
\item
  If \(k = \mathbf{R}\) or \(\mathbf{C}\) , \(\operatorname{Aut}_e(\mathfrak{g})\) is the identity component of \(\operatorname{Aut}(\mathfrak{g})\) , in other words \(\operatorname{Int}(\mathfrak{g})\) .
\end{enumerate}

By Chap. III, §3, no. 8, Cor. 2 of Prop. 29, and no. 10, Prop. 36, \(\operatorname{Aut}(\mathfrak{g}, \mathfrak{h})\) is a Lie subgroup of \(\operatorname{Aut}(\mathfrak{g})\) whose Lie algebra is the set of \(\operatorname{ad} x (x \in \mathfrak{g})\) such that \((\operatorname{ad} x) \mathfrak{h} \subset \mathfrak{h}\) , in other words \(\operatorname{ad} \mathfrak{h}\) .

Let \(\mathrm{H} \in \mathfrak{h}\) . There exists \(\varepsilon > 0\) with the following properties: for \(t \in k\) and \(|t| < \varepsilon\) , \(\exp(t\gamma(\mathrm{H}))\) is defined for all \(\gamma \in \mathrm{P}(\mathrm{R})\) , and the map \(\gamma \mapsto \exp(t\gamma(\mathrm{H}))\) is a homomorphism \(\sigma_t\) from \(\mathrm{P}(\mathrm{R})\) to \(k^*\) . For \(|t| < \varepsilon\) , \(\exp(t\operatorname{ad}\mathrm{H})\) is defined, induces the identity on \(\mathfrak{h}\) and induces on \(\mathfrak{g}^\alpha\) the homothety with ratio \(\sigma_t(\alpha)\) ; hence \(\exp t\operatorname{ad}\mathrm{H} \in (q \circ f)(\mathrm{T}_{\mathrm{P}})\) . This proves, in view of (i), that \((q \circ f)(\mathrm{T}_{\mathrm{P}})\) contains a neighbourhood of 1 in \(\operatorname{Aut}(\mathfrak{g},\mathfrak{h})\) , and consequently is an open subgroup of \(\operatorname{Aut}(\mathfrak{g},\mathfrak{h})\) . A fortiori, \(f(\mathrm{T}_{\mathrm{Q}})\) is an open subgroup of \(\operatorname{Aut}(\mathfrak{g},\mathfrak{h})\) .

For all \(\alpha \in \mathrm{R}\) , \(\exp \operatorname{ad} \mathfrak{g}^{\alpha} \subset \operatorname{Aut}_{e}(\mathfrak{g})\) . In view of (ii), \(\operatorname{Aut}_{e}(\mathfrak{g})\) contains a neighbourhood of 1 in \(\operatorname{Aut}(\mathfrak{g})\) , which proves (iii).

Assume that \(k = \mathbf{R}\) or \(\mathbf{C}\) . Then \(\mathrm{Aut}_e(\mathfrak{g})\) is contained in the identity component \(\mathrm{C}\) of \(\mathrm{Aut}(\mathfrak{g})\) (Chap. VII, §3, no. 1), and is open in \(\mathrm{Aut}(\mathfrak{g})\) by (iii). Thus \(\mathrm{Aut}_e(\mathfrak{g}) = \mathrm{C}\) . Finally, \(\mathrm{C} = \mathrm{Int}(\mathfrak{g})\) by Chap. III, §9, no. 8, Prop. 30 (i).

\section{MODULES OVER A SPLIT SEMI-SIMPLE LIE ALGEBRA}

In this paragraph, \((\mathfrak{g},\mathfrak{h})\) denotes a split semi-simple Lie algebra, R its root system, W its Weyl group, B a basis of R, \(R_{+}\) (resp. \(R_{-}\) ) the set of positive (resp. negative) roots relative to B. Put

\[
\mathfrak {n} _ {+} = \sum_ {\alpha \in R _ {+}} \mathfrak {g} ^ {\alpha}, \quad \mathfrak {n} _ {-} = \sum_ {\alpha \in R _ {-}} \mathfrak {g} ^ {\alpha}, \quad \mathfrak {b} _ {+} = \mathfrak {h} + \mathfrak {n} _ {+} \text {   and   } \mathfrak {b} _ {-} = \mathfrak {h} + \mathfrak {n} _ {-}.
\]

We have \(\mathfrak{n}_{+} = [\mathfrak{b}_{+},\mathfrak{b}_{+}],\mathfrak{n}_{-} = [\mathfrak{b}_{-},\mathfrak{b}_{-}].\)

For all \(\alpha\in\mathrm{R}\) , choose an element \(X_{\alpha}\in\mathfrak{g}^{\alpha}\) such that

\[
[ X _ {\alpha}, X _ {- \alpha} ] = - H _ {\alpha}
\]

(§2, no. 4); none of the definitions below will depend on this choice.

\subsection{WEIGHTS AND PRIMITIVE ELEMENTS}

Let V be a g-module. For all \(\lambda\in h^{*}\) , denote by \(V^{\lambda}\) the primary subspace, relative to \(\lambda\) , of V considered as an h-module (Chap. VII, §1, no. 1). The elements of \(V^{\lambda}\) are called the elements of weight \(\lambda\) of the g-module V. The sum of the \(V^{\lambda}\) is direct (Chap. VII, §1, no. 1, Prop. 3). For all \(\alpha\in h^{*}\) and \(\lambda\in h^{*}\) , \(g^{\alpha}V^{\lambda}\subset V^{\alpha+\lambda}\) (Chap. VII, §1, no. 3, Prop. 10 (ii)). The dimension of \(V^{\lambda}\) is called the multiplicity of \(\lambda\) in V; if it is \(\geq1\) , i.e.~if \(V^{\lambda}\neq0\) , \(\lambda\) is said to be a weight of V. If V is finite dimensional, the homotheties of V defined by the elements of h are semi-simple, so \(V^{\lambda}\) is the set of \(x\in V\) such that \(Hx=\lambda(H)x\) for all \(H\in h\) .

Lemma 1. Let \(\mathrm{V}\) be a \(\mathfrak{g}\) -module and \(v \in \mathrm{V}\) . The following conditions are equivalent:

\begin{enumerate}
\def\labelenumi{(\roman{enumi})}
\item
  \(\mathfrak{b}_{+}v\subset kv;\)
\item
  \(\mathfrak{h}v \subset kv\) and \(n_{+}v = 0;\)
\item
  \(\mathfrak{h}v\subset kv\) and \(\mathfrak{g}^{\alpha}v = 0\) for all \(\alpha \in \mathrm{B}\) .
\item
  \(\Longrightarrow\) (ii): Assume that \(\mathfrak{b}_{+}v\subset kv\) . There exists \(\lambda \in \mathfrak{h}^*\) such that \(v\in V^{\lambda}\) . Let \(\alpha \in \mathrm{R}_{+}\) . Then \(\mathfrak{g}^{\alpha}.v\subset V^{\lambda}\cap V^{\lambda +\alpha} = 0\) . Hence \(\mathfrak{n}_{+}v = 0\) .
\item
  \(\Longrightarrow\) (iii): This is clear.
\item
  \(\Longrightarrow\) (i): This follows from the fact that \((X_{\alpha})_{\alpha \in \mathrm{B}}\) generates \(\mathfrak{n}_{+}\) (§3, no. 3, Prop. 9 (iii)).
\end{enumerate}

DEFINITION 1. Let V be a g-module and \(v \in V\) . Then v is said to be a primitive element of V if \(v \neq 0\) and v satisfies the conditions of Lemma 1.

A primitive element belongs to one of the \(V^{\lambda}\) . For all \(\lambda\in h^{*}\) , \(V_{\pi}^{\lambda}\) denotes the set of \(v\in V^{\lambda}\) such that \(b_{+}v\subset kv\) . Thus, the primitive elements of weight \(\lambda\) are the non-zero elements of \(V_{\pi}^{\lambda}\) .

PROPOSITION 1. Let V be a g-module, v a primitive element of V and \(\omega\) the weight of v. Assume that V is generated by v as a g-module.

\begin{enumerate}
\def\labelenumi{(\roman{enumi})}
\item
  If \(\mathrm{U}(\mathfrak{n}_{-})\) denotes the enveloping algebra of \(\mathfrak{n}_{-}\) , we have \(\mathrm{V} = \mathrm{U}(\mathfrak{n}_{-}).v\) .
\item
  For all \(\lambda \in \mathfrak{h}^*\) , \(V^{\lambda}\) is the set of \(x \in V\) such that \(Hx = \lambda(H)x\) for all \(H \in \mathfrak{h}\) . We have \(V = \bigoplus_{\lambda \in \mathfrak{h}^*} V^{\lambda}\) , and each \(V^{\lambda}\) is finite dimensional. The space \(V^{\omega}\) is of dimension 1, and every weight of \(V\) is of the form \(\omega - \sum_{\alpha \in B} n_{\alpha}. \alpha\) , where the \(n_{\alpha}\) are integers \(\geq 0\) .
\item
  V is an indecomposable \(\mathfrak{g}\) -module, and its commutant reduces to the scalars.
\item
  Let \(\mathrm{U}(\mathfrak{g})\) be the enveloping algebra of g, and Z the centre of \(\mathrm{U}(\mathfrak{g})\) . There exists a unique homomorphism \(\chi\) from Z to k such that, for all \(z \in Z\) , \(z_{V}\) is the homothety with ratio \(\chi(z)\) .
\end{enumerate}

Let \(\mathrm{U}(\mathfrak{b}_{+})\) be the enveloping algebra of \(\mathfrak{b}_{+}\) . We have \(\mathrm{U}(\mathfrak{g}) = \mathrm{U}(\mathfrak{n}_{-})\) . \(\mathrm{U}(\mathfrak{b}_{+})\) (Chap. I, §2, no. 7, Cor. 6 of Th. 1). Hence

\[
\mathrm{V} = \mathrm{U} (\mathfrak {g}). v = \mathrm{U} (\mathfrak {n} _ {-}). \mathrm{U} (\mathfrak {b} _ {+}). v = \mathrm{U} (\mathfrak {n} _ {-}). v.
\]

Denote by \(\alpha_{1},\ldots,\alpha_{n}\) the distinct elements of \(R_{+}\) . Then

\[
\left(X _ {- \alpha_ {1}} ^ {p _ {1}} X _ {- \alpha_ {2}} ^ {p _ {2}} \dots X _ {- \alpha_ {n}} ^ {p _ {n}}\right) (p _ {1}, \ldots , p _ {n}) \in \mathbf {N} ^ {n}
\]

is a basis of \(\mathrm{U}(\mathfrak{n}_{-})\) , so

\[
\mathrm{V} = \sum_ {(p _ {1}, \dots , p _ {n}) \in \mathbf {N} ^ {n}} k X _ {- \alpha_ {1}} ^ {p _ {1}} \dots X _ {- \alpha_ {n}} ^ {p _ {n}} v.\tag{1}
\]

For \(\lambda \in \mathfrak{h}^*\) , put

\[
\mathrm{T} _ {\lambda} = \sum_ {(p _ {1}, \ldots , p _ {n}) \in \mathbf {N} ^ {n}, \omega - p _ {1} \alpha_ {1} - \dots - p _ {n} \alpha_ {n} = \lambda} k X _ {- \alpha_ {1}} ^ {p _ {1}} \ldots X _ {- \alpha_ {n}} ^ {p _ {n}} v.
\]

By Chap. VII, §1, no. 1, Prop. 2 (ii), if \(h \in \mathfrak{h}\) , \(h_{\mathrm{V}}|_{\mathrm{T}_{\lambda}}\) is the homothety with ratio \(\lambda(h)\) . So \(\mathrm{T}_{\lambda} \subset \mathrm{V}^{\lambda}\) . On the other hand, (1) implies that

\[
\mathrm{V} = \sum_ {\lambda \in \omega - \mathbf {N} \alpha_ {1} - \dots - \mathbf {N} \alpha_ {n}} \mathrm{T} _ {\lambda}.
\]

The sum of the \(V^{\lambda}\) is direct (Chap. VII, §1, no. 1, Prop. 3). From these observations it follows that \(V^{\lambda} = T_{\lambda}\) , that V is the direct sum of the \(V^{\lambda}\) , and that \(V^{\lambda}\) is the set of \(x \in V\) such that \(hx = \lambda(h)x\) for all \(h \in h\) . On the other hand, \(\dim V^{\lambda}\) is at most the cardinal of the set of \((p_{1}, \ldots, p_{n}) \in \mathbf{N}^{n}\) such that \(p_{1}\alpha_{1} + \cdots + p_{n}\alpha_{n} = \omega - \lambda\) . This proves that \(V^{\lambda} = 0\) if \(\omega - \lambda \notin \sum_{\alpha \in B} N\alpha\) , that \(\dim V^{\omega} = 1\) , and that the \(V^{\lambda}\) are all finite dimensional.

Let \(c\) be an element of the commutant of V. For all \(h \in \mathfrak{h}\) ,

\[
h c (v) = c h (v) = \omega (h) c (v),
\]

so \(c(v) \in V^{\omega}\) ; hence there exists \(t \in k\) such that \(c(v) = tv\) . Now, for all \((p_1, \ldots, p_n) \in \mathbf{N}^n\) ,

\[
c X _ {- \alpha_ {1}} ^ {p _ {1}} \dots X _ {- \alpha_ {n}} ^ {p _ {n}} v = X _ {- \alpha_ {1}} ^ {p _ {1}} \dots X _ {- \alpha_ {n}} ^ {p _ {n}} c v = t X _ {- \alpha_ {1}} ^ {p _ {1}} \dots X _ {- \alpha_ {n}} ^ {p _ {n}} v
\]

so that \(c = t.1\) . Hence, the commutant of \(V\) reduces to the scalars. This implies (iv) and the fact that \(V\) is indecomposable.

DEFINITION 2. The homomorphism \(\chi\) of Prop. 1 (iv) is called the central character of the \(\mathfrak{g}\) -module V.

PROPOSITION 2. Let V be a g-module generated by a primitive element e of weight \(\omega\) , and X a semi-simple g-module. Let \(\Phi\) be the set of homomorphisms from the g-module V to the g-module X. Then \(\varphi\mapsto\varphi(e)\) is an isomorphism from the vector space \(\Phi\) to the vector space \(X_{\pi}^{\omega}\) .

It is clear that \(\varphi(e) \in X_{\pi}^{\omega}\) for all \(\varphi \in \Phi\) . If \(\varphi \in \Phi\) and \(\varphi(e) = 0\) , then \(\varphi = 0\) since \(e\) generates the \(\mathfrak{g}\) -module V. We show that, if \(f\) is a non-zero element of \(X_{\pi}^{\omega}\) , there exists \(\varphi \in \Phi\) such that \(\varphi(e) = f\) . Let \(X'\) be the submodule of X generated by \(f\) . By Prop. 1, \(X'\) is indecomposable, hence simple since X is semi-simple. The element \((e, f)\) is primitive in the \(\mathfrak{g}\) -module \(V \times X\) . Let N be the submodule of \(V \times X\) generated by \((e, f)\) . Then \(N \cap X \subset \mathrm{pr}_2(N) = X'\) , so \(N \cap X = 0\) or \(X'\) ; if \(N \cap X = X'\) , N contains the linearly independent elements \((e, f)\) and \((0, f)\) which are primitive of weight \(\omega\) ; this is absurd (Prop. 1), so \(N \cap X = 0\) . Thus \(\mathrm{pr}_1|_{N}\) is an injective map \(h\) from N to V; this map is surjective since its image contains \(e\) . Thus \(\varphi = \mathrm{pr}_2 \circ h^{-1}\) is a homomorphism from the \(\mathfrak{g}\) -module V to the \(\mathfrak{g}\) -module X such that \(\varphi(e) = f\) .

\subsection{SIMPLE MODULES WITH A HIGHEST WEIGHT}

Recall that fixing B defines an order relation on \(h_{Q}^{*}\) (Chap. VI, §1, no. 6). The elements of \(h_{Q}^{*}\) that are \(\geq 0\) are the linear combinations of elements of B with rational coefficients \(\geq 0\) .

More generally, we shall consider the following order relation between elements \(\lambda, \mu \in \mathfrak{h}^*\) :

\(\lambda - \mu\) is a linear combination of elements of B with rational coefficients \(\geq 0\) .

Lemma 2. Let V be a simple g-module, \(\omega\) a weight of V. The following conditions are equivalent:

\begin{enumerate}
\def\labelenumi{(\roman{enumi})}
\item
  every weight of V is of the form \(\omega -\mu\) where \(\mu\) is a radical weight \(\geq 0\) ;
\item
  \(\omega\) is the highest weight of V;
\item
  for all \(\alpha \in \mathrm{B}\) , \(\omega + \alpha\) is not a weight of V;
\item
  there exists a primitive element of weight \(\omega\) .
\item
  \(\Longrightarrow\) (ii) \(\Longrightarrow\) (iii): This is clear.
\item
  \(\Longrightarrow\) (iv): Assume that condition (iii) is satisfied. For all \(h\in \mathfrak{h}\) ,
\end{enumerate}

\(\operatorname {Ker}(h_{\mathrm{V}} - \omega (h))\)

is non-zero, contained in \(V^{\omega}\) , and stable under \(h_{V}\) . By induction on dim h, we see that there exists a non-zero v in \(V^{\omega}\) such that \(h v \subset kv\) . Condition (iii) implies that \(n_{+}v = 0\) , so v is primitive.

\begin{enumerate}
\def\labelenumi{(\roman{enumi})}
\setcounter{enumi}{3}
\tightlist
\item
  \(\Longrightarrow\) (i): Let \(v\) be a primitive element of weight \(\omega\) . Since V is simple, V is generated by \(v\) as a \(\mathfrak{g}\) -module. Assertion (i) now follows from Prop. 1.
\end{enumerate}

Q.E.D.

Thus, for any simple g-module, the existence of a primitive element is equivalent to that of a highest weight, or to that of a maximal weight.

There exist simple \(\mathfrak{sl}(2,\mathbf{C})\) -modules V that have no weights for any Cartan subalgebra \(\mathfrak{h}\) of \(\mathfrak{sl}(2,\mathbf{C})\) (§1, Exerc. 14 \(f\) )). These modules are of infinite dimension over \(\mathbf{C}\) (§1, no. 3, Th. 1).

PROPOSITION 3. Let V be a simple \(\mathfrak{g}\) -module with a highest weight \(\omega\) .

\begin{enumerate}
\def\labelenumi{(\roman{enumi})}
\item
  The primitive elements of V are the non-zero elements of \(V^{\omega}\) .
\item
  V is semi-simple as an \(\mathfrak{h}\) -module.
\item
  We have \(V = \bigoplus_{\lambda \in h^{*}} V^{\lambda}\) . For all \(\lambda \in h^{*}\) , \(V^{\lambda}\) is finite dimensional. We have \(\dim V^{\omega} = 1\) .
\item
  The \(\mathfrak{g}\) -module V is absolutely simple.
\end{enumerate}

Assertions (i), (ii) and (iii) follow from Prop. 1 and Lemma 2. Assertion (iv) follows from Prop. 1 (iii) and Algebra, Chap. VIII, §7, no. 3.

COROLLARY. If \(\mathrm{V}\) is finite dimensional, the canonical homomorphism \(\mathrm{U}(\mathfrak{g}) \to \operatorname{End}(\mathrm{V})\) is surjective.

This follows from (iv), cf.~Algebra, Chap. VIII, §3, no. 3.

PROPOSITION 4. Let V be a simple g-module with a highest weight \(\omega\) , X a semi-simple g-module, and \(X'\) the isotypical component of type V in X. Then \(X'\) is the submodule of X generated by \(X_{\pi}^{\omega}\) . Its length is equal to the dimension of \(X_{\pi}^{\omega}\) .

Let \(X''\) be the submodule of X generated by \(X_{\pi}^{\omega}\) . It is clear that every submodule of X isomorphic to V is contained in \(X''\) . Hence \(X' \subset X''\) . On the other hand, let \(\varPhi\) be the set of homomorphisms from the g-module V to the g-module X. The length of \(X'\) is \(\dim_k \varPhi\) (Algebra, Chap. VIII, §4, no. 4), that is \(\dim_k X_{\pi}^{\omega}\) (Prop. 2).

\subsection{EXISTENCE AND UNIQUENESS THEOREM}

Let \(\lambda \in \mathfrak{h}^*\) . Since \(\mathfrak{b}_{+} = \mathfrak{h} \oplus \mathfrak{n}_{+}\) and since \(\mathfrak{n}_{+} = [\mathfrak{b}_{+}, \mathfrak{b}_{+}]\) , the map \(h + n \mapsto \lambda(h)\) (where \(h \in \mathfrak{h}\) , \(n \in \mathfrak{n}_{+}\) ) from \(\mathfrak{b}_{+}\) to \(k\) is a 1-dimensional representation of \(\mathfrak{b}_{+}\) . Denote by \(\mathrm{L}_{\lambda}\) the \(k\) -vector space \(k\) equipped with the \(\mathfrak{b}_{+}\) -module structure defined by this representation. Let \(\mathrm{U}(\mathfrak{g})\) , \(\mathrm{U}(\mathfrak{b}_{+})\) be the enveloping algebras of \(\mathfrak{g}\) , \(\mathfrak{b}_{+}\) , so that \(\mathrm{U}(\mathfrak{b}_{+})\) is a subalgebra of \(\mathrm{U}(\mathfrak{g})\) ; recall that \(\mathrm{U}(\mathfrak{g})\) is a free right \(\mathrm{U}(\mathfrak{b}_{+})\) -module (Chap. I, §2, no. 7, Cor. 5 of Th. 1). Put

\[
\mathrm{Z} (\lambda) = \mathrm{U} (\mathfrak {g}) \otimes_ {\mathrm{U} (\mathfrak {b} _ {+})} \mathrm{L} _ {\lambda}.\tag{2}
\]

Then \(Z(\lambda)\) is a left \(\mathfrak{g}\) -module. Denote by \(e\) the element \(1 \otimes 1\) of \(Z(\lambda)\) .

PROPOSITION 5. (i) The element \(e\) of \(Z(\lambda)\) is primitive of weight \(\lambda\) and generates the \(\mathfrak{g}\) -module \(Z(\lambda)\) .

\begin{enumerate}
\def\labelenumi{(\roman{enumi})}
\setcounter{enumi}{1}
\item
  Let \(\mathrm{Z}^{+}(\lambda) = \sum_{\lambda \neq \mu} \mathrm{Z}(\lambda)^{\mu}\) . Every submodule of \(\mathrm{Z}(\lambda)\) distinct from \(\mathrm{Z}(\lambda)\) is contained in \(\mathrm{Z}^{+}(\lambda)\) .
\item
  There exists a largest submodule \(F_{\lambda}\) of \(Z(\lambda)\) distinct from \(Z(\lambda)\) . The quotient module \(Z(\lambda)/F_{\lambda}\) is simple and has highest weight \(\lambda\) .
\end{enumerate}

It is clear that e generates the g-module \(Z(\lambda)\) . If \(x \in b_{+}\) ,

\[
x. e = (x. 1) \otimes 1 = (1. x) \otimes 1 = 1 \otimes x. 1 = \lambda (x) (1 \otimes 1) = \lambda (x) e,
\]

hence (i).

The \(\mathfrak{h}\) -module \(Z(\lambda)\) is semi-simple (Prop. 1). If \(G\) is a \(\mathfrak{g}\) -submodule of \(Z(\lambda)\) , then \(G = \sum_{\mu \in \mathfrak{h}^*}(G \cap Z(\lambda)^\mu)\) . The hypothesis \(G \cap Z(\lambda)^\lambda \neq 0\) implies that \(G = Z(\lambda)\) , since \(\dim Z(\lambda)^\lambda = 1\) and \(e\) generates the \(\mathfrak{g}\) -module \(Z(\lambda)\) . If \(G \neq Z(\lambda)\) , then \(G = \sum_{\mu \neq \lambda} G \cap Z(\lambda)^\mu \subset Z^+( \lambda)\) .

Let \(F_{\lambda}\) be the sum of the g-submodules of \(Z(\lambda)\) distinct from \(Z(\lambda)\) . By (ii), \(F_{\lambda} \subset Z^{+}(\lambda)\) . Hence \(F_{\lambda}\) is the largest submodule of \(Z(\lambda)\) distinct from \(Z(\lambda)\) . It is clear that \(Z(\lambda)/F_{\lambda}\) is simple and that the canonical image of e in \(Z(\lambda)/F_{\lambda}\) is primitive of weight \(\lambda\) .

In the remainder of this chapter, the g-module \(Z(\lambda)/F_{\lambda}\) of Prop. 5 will be denoted by \(E(\lambda)\) .

THEOREM 1. Let \(\lambda \in \mathfrak{h}^*\) . The \(\mathfrak{g}\) -module \(\mathrm{E}(\lambda)\) is simple and has highest weight \(\lambda\) . Every simple \(\mathfrak{g}\) -module of highest weight \(\lambda\) is isomorphic to \(\mathrm{E}(\lambda)\) .

The first assertion follows from Prop. 5 (iii). The second follows from Prop. 4.

PROPOSITION 6. Let V be a g-module, \(\lambda\) an element of \(h^{*}\) and v a primitive element of V of weight \(\lambda\) .

\begin{enumerate}
\def\labelenumi{(\roman{enumi})}
\item
  There exists a unique homomorphism of \(\mathfrak{g}\) -modules \(\psi: \mathrm{Z}(\lambda) \to \mathrm{V}\) such that \(\psi(e) = v\) .
\item
  Assume that v generates V. Then \(\psi\) is surjective. Moreover, \(\psi\) is bijective if and only if, for every non-zero element u of \(\mathrm{U}(\mathfrak{n}_{-})\) , \(u_{V}\) is injective.
\item
  The map \(u \mapsto u \otimes 1\) from \(\mathrm{U}(\mathfrak{n}_{-})\) to \(\mathrm{Z}(\lambda)\) is bijective.
\end{enumerate}

Let K be the kernel of the representation of \(\mathrm{U}(\mathfrak{b}_{+})\) on \(L_{\lambda}\) ; it is of codimension 1 in \(\mathrm{U}(\mathfrak{b}_{+})\) . Let \(\mathrm{J} = \mathrm{U}(\mathfrak{g})\mathrm{K}\) be the left ideal of \(\mathrm{U}(\mathfrak{g})\) generated by K; then \(L_{\lambda}\) can be identified with \(\mathrm{U}(\mathfrak{b}_{+})/\mathrm{K}\) as a left \(\mathrm{U}(\mathfrak{b}_{+})\) -module, and \(\mathrm{Z}(\lambda)\) can be identified with \(\mathrm{U}(\mathfrak{g})/\mathrm{J}\) as a left \(\mathrm{U}(\mathfrak{g})\) -module. We have K.v = 0, so J.v = 0, which proves (i).

Now assume that \(v\) generates V. It is clear that \(\psi\) is surjective.

By the Poincaré-Birkhoff-Witt theorem (Chap. I, §2, no. 7, Cor. 6 of Th. 1), a basis of \(\mathrm{U}(\mathfrak{n}_{-})\) over \(k\) is also a basis of \(\mathrm{U}(\mathfrak{g})\) as a right \(\mathrm{U}(\mathfrak{b}_{+})\) -module. Hence the map \(\varphi : u \mapsto u \otimes 1\) from \(\mathrm{U}(\mathfrak{n}_{-})\) to \(\mathrm{U}(\mathfrak{g}) \otimes_{\mathrm{U}(\mathfrak{b}_{+})} \mathrm{L}_{\lambda}\) is bijective. Let \(u \in \mathrm{U}(\mathfrak{n}_{-})\) . Then \(\varphi^{-1} \circ u_{\mathrm{Z}(\lambda)} \circ \varphi\) is left multiplication by \(u\) on \(\mathrm{U}(\mathfrak{n}_{-})\) . In view of Chap. I, §2, no. 7, Cor. 7 of Th. 1, \(u_{\mathrm{Z}(\lambda)}\) is injective if \(u \neq 0\) . Consequently, if \(\psi\) is bijective, then \(u_{\mathrm{V}}\) is injective for non-zero \(u\) in \(\mathrm{U}(\mathfrak{n}_{-})\) .

Assume that \(\psi\) is not injective. There exists \(u \in \mathrm{U}(\mathfrak{n}_{-})\) such that \(u \neq 0\) and \(\psi(\varphi(u)) = 0\) . Then

\[
u _ {\mathrm{V}}. v = u _ {\mathrm{V}}. \psi (1 \otimes 1) = \psi (u \otimes 1) = \psi (\varphi (u)) = 0.
\]

COROLLARY 1. Let \(\lambda \in \mathfrak{h}^*\) and \(\alpha \in B\) be such that \(\lambda(H_{\alpha}) + 1 \in \mathbf{N}\) . Then \(Z(-\alpha + s_{\alpha}\lambda)\) is isomorphic to a \(\mathfrak{g}\) -submodule of \(Z(\lambda)\) .

Put \(m = \lambda(H_{\alpha})\) . Let \(x = X_{-\alpha}^{m+1}.e \in Z(\lambda)\) , and let V be the submodule of \(Z(\lambda)\) generated by x. Then \(x \neq 0\) (Prop. 6). On the other hand, \(x \in Z(\lambda)^{\lambda-(m+1)\alpha}\) . For \(\beta \in B\) and \(\beta \neq \alpha\) , \([\mathfrak{g}^{-\alpha}, \mathfrak{g}^{\beta}] = 0\) and \(g^{\beta}.e = 0\) , so \(g^{\beta}.x = 0\) . Finally, since \([X_{\alpha}, X_{-\alpha}] = -H_{\alpha}\) , we have

\[
[ X _ {\alpha}, X _ {- \alpha} ^ {m + 1} ] = (m + 1) X _ {- \alpha} ^ {m} (- H _ {\alpha} + m)
\]

(§1, no. 1, Lemma 1 (ii)), so

\[
X _ {\alpha}. x = X _ {\alpha} X _ {- \alpha} ^ {m + 1}. e = [ X _ {\alpha}, X _ {- \alpha} ^ {m + 1} ]. e = (m + 1) X _ {- \alpha} ^ {m} (m e - \lambda (H _ {\alpha}) e) = 0.
\]

Thus, \(x\) is primitive of weight \(\lambda - (m + 1)\alpha\) . In view of Prop. 6, the \(\mathfrak{g}\) -module \(V\) is isomorphic to \(Z(-\alpha + \lambda - m\alpha) = Z(-\alpha + s_{\alpha}\lambda)\) .

COROLLARY 2. Let \(\rho = \frac{1}{2}\sum_{\alpha \in \mathrm{R}_{+}}\alpha\) , and \(\lambda, \mu \in \mathfrak{h}^*\) . Assume that \(\lambda + \rho\) is a dominant weight in \(\mathrm{R}\) , and that there exists \(w \in \mathrm{W}\) with \(\mu + \rho = w(\lambda + \rho)\) . Then \(\mathrm{Z}(\mu)\) is isomorphic to a submodule of \(\mathrm{Z}(\lambda)\) .

The assertion is clear when w = 1. Assume that it is established whenever w is of length \textless{} q. If w is of length q, there exists \(\alpha \in B\) such that \(w = s_{\alpha}w'^{-1}\) , with \(l(w') = q - 1\) . We have \(w'(\alpha) \in \mathrm{R}_{+}\) (Chap. VI, §1, no. 6, Cor. 2 of Prop. 17), and hence \(w'^{-1}(\lambda + \rho)(H_{\alpha}) = (\lambda + \rho)(H_{w'\alpha})\) is an integer \(\geq 0\) . Put

\[
\mu^ {\prime} = w ^ {- 1} (\lambda + \rho) - \rho .
\]

By the induction hypothesis, \(Z(\mu')\) is isomorphic to a submodule of \(Z(\lambda)\) . On the other hand, by Chap. VI, §1, no. 10, Prop. 29 (ii),

\[
- \alpha + s _ {\alpha} \mu^ {\prime} = - \alpha + s _ {\alpha} w ^ {- 1} (\lambda + \rho) - s _ {\alpha} \rho = w (\lambda + \rho) - \rho = \mu .
\]

Moreover, \(\rho(H_{\alpha}) = 1\) (Chap. VI, §1, Prop. 29 (iii)), so \(\mu'(H_{\alpha}) + 1 \in \mathbf{N}\) . Cor. 1 now implies that \(Z(\mu)\) is isomorphic to a submodule of \(Z(\mu')\) , and hence also to a submodule of \(Z(\lambda)\) .

\subsection{\texorpdfstring{COMMUTANT OF \(\mathfrak{h}\) IN THE ENVELOPING ALGEBRA OF \(\mathfrak{g}\)}{COMMUTANT OF \textbackslash mathfrak\{h\} IN THE ENVELOPING ALGEBRA OF \textbackslash mathfrak\{g\}}}

Let U be the enveloping algebra of g, \(V \subset U\) the enveloping algebra of h. The algebra V can be identified with the symmetric algebra \(\mathbf{S}(\mathfrak{h})\) of h, and also with the algebra of polynomial functions on \(h^{*}\) . Denote by \(\alpha_{1},\ldots,\alpha_{n}\) the pairwise distinct positive roots. Let \((H_{1},\ldots,H_{l})\) be a basis of h. By the Poincaré-Birkhoff-Witt theorem, the elements

\[
u ((q _ {i}), (m _ {i}), (p _ {i})) = X _ {- \alpha_ {1}} ^ {q _ {1}} \dots X _ {- \alpha_ {n}} ^ {q _ {n}} H _ {1} ^ {m _ {1}} \dots H _ {l} ^ {m _ {l}} X _ {\alpha_ {1}} ^ {p _ {1}} \dots X _ {\alpha_ {n}} ^ {p _ {n}}
\]

\((q_{i},m_{i},p_{i})\) integers \(\geq 0\) ) form a basis of the vector space U. For all \(h\in \mathfrak{h}\) , we have

\[
[ h, u ((q _ {i}), (m _ {i}), (p _ {i})) ] = ((p _ {1} - q _ {1}) \alpha_ {1} + \dots + (p _ {n} - q _ {n}) \alpha_ {n}) (h) u ((q _ {i}), (m _ {i}), (p _ {i})).\tag{3}
\]

The vector space U is a g-module (hence also an h-module) under the adjoint representation. If \(\lambda \in h^{*}\) , the subspaces \(U^{\lambda}\) and \(U_{\lambda}\) are defined (Chap. VII, §1, no. 3); formula (3) shows that \(U^{\lambda} = U_{\lambda}\) and that \(U = \bigoplus_{\lambda \in Q} U^{\lambda}\) (where Q is the group of radical weights of R). In particular, \(U^{0}\) is the commutant of h, or of V, in U.

Lemma 3. Put \(\mathrm{L} = (\mathfrak{n}_{-}\mathrm{U})\cap \mathrm{U}^{0}\)

\begin{enumerate}
\def\labelenumi{(\roman{enumi})}
\item
  We have \(\mathrm{L} = (\mathrm{U}\mathfrak{n}_{+})\cap \mathrm{U}^{0}\) , and \(\mathrm{L}\) is a two-sided ideal of \(\mathrm{U}^0\) .
\item
  We have \(U^{0} = V \oplus L\) .
\end{enumerate}

It is clear that \(\mathfrak{n}_{-}\mathrm{U}\) (resp. \(\mathrm{U}\mathfrak{n}_{+}\) ) is the set of linear combinations of the elements \(u((q_{i}),(m_{i}),(p_{i}))\) such that \(\sum q_{i}>0\) (resp. \(\sum p_{i}>0\) ). On the other hand

\[
u \left(\left(q _ {i}\right), \left(m _ {i}\right), \left(p _ {i}\right)\right) \in \mathrm{U} ^ {0} \Longleftrightarrow p _ {1} \alpha_ {1} + \dots + p _ {n} \alpha_ {n} = q _ {1} \alpha_ {1} + \dots + q _ {n} \alpha_ {n}.
\]

This implies that \((\mathfrak{n}_{-}\mathrm{U})\cap\mathrm{U}^{0}=(\mathrm{U}\mathfrak{n}_{+})\cap\mathrm{U}^{0}\) . Finally, \((\mathfrak{n}_{-}\mathrm{U})\cap\mathrm{U}^{0}\) (resp. \((\mathrm{U}\mathfrak{n}_{+})\cap\mathrm{U}^{0}\) ) is a right (resp. left) ideal of \(U^{0}\) , hence (i). Further, an element \(u((q_{i}),(m_{i}),(p_{i}))\) that is in \(U^{0}\) belongs to V (resp. to L) if and only if \(p_{1}=\cdots=p_{n}=q_{1}=\cdots=q_{n}=0\) (resp. \(p_{1}+\cdots+p_{n}+q_{1}+\cdots+q_{n}>0\) ), hence (ii). Q.E.D.

In view of Lemma 3, the projection of \(U^{0}\) onto V with kernel L is a homomorphism of algebras. It is called the Harish-Chandra homomorphism from \(U^{0}\) to V (relative to B). Recall that V can be identified with the algebra of polynomial functions on \(h^{*}\) .

PROPOSITION 7. Let \(\lambda \in \mathfrak{h}^*\) , E a \(\mathfrak{g}\) -module generated by a primitive element of weight \(\lambda\) , \(\chi\) the central character of E, and \(\varphi\) the Harish-Chandra homomorphism from \(\mathrm{U}^0\) to V. Then, \(\chi(z) = (\varphi(z))(\lambda)\) for all \(z\) in the centre of U.

Let v be a primitive element of E of weight \(\lambda\) , and z an element of the centre of U. There exist \(u_{1},\ldots,u_{p}\in U\) and \(n_{1},\ldots,n_{p}\in\mathfrak{n}_{+}\) such that \(z=\varphi(z)+u_{1}n_{1}+\cdots+u_{p}n_{p}\) . Then

\[
\chi (z) v = z v = \varphi (z) v + u _ {1} n _ {1} v + \dots + u _ {p} n _ {p} v = \varphi (z) v = (\varphi (z)) (\lambda) v.
\]

COROLLARY. Let \(\langle\cdot,\cdot\rangle\) be a non-degenerate invariant symmetric bilinear form on g, C the Casimir element associated to \(\langle\cdot,\cdot\rangle\) . Denote also by \(\langle\cdot,\cdot\rangle\) the inverse form on \(h^{*}\) of the restriction of \(\langle\cdot,\cdot\rangle\) to h (§2, no. 3, Prop. 5). Then \(\chi(\mathrm{C})=\langle\lambda,\lambda+2\rho\rangle\) , where \(\rho=\frac{1}{2}\sum_{\alpha\in R_{+}}\alpha\) .

We recall the notations of §2, no. 3, Prop. 6. We have

\[
\begin{array}{r} \mathrm{C} = \sum_ {\alpha \in \mathrm{R} _ {-}} \frac {1}{\langle X _ {\alpha} , X _ {- \alpha} \rangle} X _ {\alpha} X _ {- \alpha} + \sum_ {\alpha \in \mathrm{R} _ {+}} \frac {1}{\langle X _ {\alpha} , X _ {- \alpha} \rangle} X _ {- \alpha} X _ {\alpha} \\ + \sum_ {\alpha \in \mathrm{R} _ {+}} \frac {1}{\langle X _ {\alpha} , X _ {- \alpha} \rangle} [ X _ {\alpha}, X _ {- \alpha} ] + \sum_ {i \in I} H _ {i} H _ {i} ^ {\prime} \end{array}
\]

so

\[
\varphi (\mathrm{C}) = \sum_ {\alpha \in \mathrm{R} _ {+}} \frac {1}{\langle X _ {\alpha} , X _ {- \alpha} \rangle} [ X _ {\alpha}, X _ {- \alpha} ] + \sum_ {i \in I} H _ {i} H _ {i} ^ {\prime}.
\]

By Prop. 7,

\[
\chi (\mathrm{C}) = \sum_ {\alpha \in \mathrm{R} _ {+}} \frac {1}{\langle X _ {\alpha} , X _ {- \alpha} \rangle} \lambda ([ X _ {\alpha}, X _ {- \alpha} ]) + \sum_ {i \in I} \lambda (H _ {i}) \lambda (H _ {i} ^ {\prime}).
\]

Let \(h_{\lambda}\) be the element of \(\mathfrak{h}\) such that \(\langle h_{\lambda}, h \rangle = \lambda(h)\) for all \(h \in \mathfrak{h}\) . By §2, no. 2, Prop. 1,

\[
\lambda \left(\frac {1}{\langle X _ {\alpha} , X _ {- \alpha} \rangle} [ X _ {\alpha}, X _ {- \alpha} ]\right) = \left\langle h _ {\lambda}, \frac {1}{\langle X _ {\alpha} , X _ {- \alpha} \rangle} [ X _ {\alpha}, X _ {- \alpha} ] \right\rangle = \alpha (h _ {\lambda}) = \langle \lambda , \alpha \rangle .
\]

Hence

\[
\chi (\mathrm{C}) = \left(\sum_ {\alpha \in \mathrm{R} _ {+}} \langle \lambda , \alpha \rangle\right) + \langle \lambda , \lambda \rangle = \langle \lambda , \lambda + 2 \rho \rangle .
\]

\section{FINITE DIMENSIONAL MODULES OVER A SPLIT SEMI-SIMPLE LIE ALGEBRA}

In this paragraph, we retain the general notations of §6. We denote by P (resp. Q) the group of weights of R (resp. radical weights of R). We denote by \(P_{+}\) (resp. \(Q_{+}\) ) the set of elements of P (resp. Q) that are positive for the order relation defined by B. We denote by \(P_{++}\) the set of dominant weights of R relative to B (Chap. VI, §1, no. 10). An element \(\lambda\) of \(h^{*}\) belongs to P (resp. to \(P_{++}\) ) if and only if all the \(\lambda(H_{\alpha})\) , \(\alpha \in B\) , are integers (resp. integers \(\geq 0\) ). We have \(P_{++} \subset P_{+}\) (Chap. VI, §1, no. 6). If \(w \in W\) , we denote by \(\varepsilon(w)\) the determinant of w, which is equal to 1 or -1. We put \(\rho = \frac{1}{2} \sum_{\alpha \in R_{+}} \alpha\) .